% Compile this file for the two-column paper, or manuscript_submission.tex
% for the single-column, double-spaced manuscript. Both share all content.
% \ifdefined\SubmissionFormat
%   % IEEEtran's 12pt unity baseline is 13.92pt; this factor gives 24pt.
%   \def\CLASSINPUTbaselinestretch{1.724137931}
%   \def\CLASSINPUTinnersidemargin{1in}
%   \def\CLASSINPUToutersidemargin{1in}
%   \def\CLASSINPUTtoptextmargin{1in}
%   \def\CLASSINPUTbottomtextmargin{1in}
\documentclass[journal,12pt,onecolumn,draftclsnofoot,letterpaper]{IEEEtran}
% \else
  % \documentclass[journal]{IEEEtran}
% \fi

\usepackage{amsmath,amssymb,amsthm,mathrsfs,mathtools,booktabs,array,microtype,graphicx}
\usepackage{tikz,pgfplots}
\usetikzlibrary{arrows.meta,positioning,calc,fit,backgrounds,patterns}
\usepgfplotslibrary{groupplots,fillbetween}
% Shared vector-figure design for the manuscript. All graphics use TikZ/PGFPlots.
\definecolor{figBlue}{HTML}{216A87}
\definecolor{figOrange}{HTML}{B75B30}
\definecolor{figTeal}{HTML}{2B8C7F}
\definecolor{figGray}{HTML}{6B7280}
\pgfplotsset{compat=1.18}
\pgfplotsset{
  paper axis/.style={
    /tikz/font=\footnotesize,
    tick label style={font=\footnotesize},
    label style={font=\footnotesize},
    title style={font=\footnotesize},
    legend style={font=\footnotesize,draw=none,fill=none},
    axis line style={black!65,line width=.45pt},
    tick style={black!65,line width=.4pt},
    major grid style={black!9,line width=.3pt},
    minor grid style={black!5,line width=.25pt},
    tick align=outside,
    scaled ticks=false,
    clip=true,
    every axis plot/.append style={line width=1pt}
  }
}
\tikzset{
  paper arrow/.style={-{Stealth[length=2mm,width=1.4mm]},line width=.75pt},
  paper block/.style={draw=figBlue,line width=.7pt,rounded corners=1.5pt,fill=figBlue!5,align=center,font=\footnotesize},
  paper note/.style={font=\footnotesize,text=figGray,align=center}
}

\usepackage{placeins,balance}
\usepackage[hidelinks]{hyperref}
\ifdefined\SubmissionFormat
  \usepackage{etoolbox,adjustbox}
  \renewcommand{\IEEEPARstart}[2]{#1#2}
  \setlength{\tabcolsep}{5pt}
  \makeatletter
  % Avoid the large publication-title gap and use a compact title block.
  \def\@IEEEstatictitlevskip{\vskip 0pt}
  \patchcmd{\@maketitle}{\Huge}{\def\baselinestretch{1}\Huge}{}%
    {\PackageError{manuscript}{Could not set title spacing}{}}
  % Keep abstract, keywords, references, and captions at the text size.
  \def\@IEEEabskeysecsize{\normalsize}
  \patchcmd{\thebibliography}{\footnotesize}{\normalsize}{}%
    {\PackageError{manuscript}{Could not set reference font size}{}}
  \let\SubmissionMakeCaption\@makecaption
  \long\def\@makecaption#1#2{\begingroup
    \let\footnotesize\normalsize
    \SubmissionMakeCaption{#1}{#2}\endgroup}
  \makeatother
  \renewenvironment{abstract}{\par%
    \noindent\textbf{Abstract}\par\noindent\ignorespaces}%
    {\par\addvspace{0.5\baselineskip}}
  \renewenvironment{IEEEkeywords}{\begin{samepage}\par\noindent
    \textbf{Index Terms---}\ignorespaces}%
    {\par\end{samepage}\addvspace{0.5\baselineskip}}
  % Preserve the original vector-figure typography and shrink only if needed.
  \newcommand{\paperfigure}[1]{\begingroup
    \def\baselinestretch{1}%
    \renewcommand{\footnotesize}{\fontsize{8}{9.5}\selectfont}%
    \renewcommand{\scriptsize}{\fontsize{7}{8}\selectfont}%
    \begin{adjustbox}{max width=\linewidth}\input{#1}\end{adjustbox}%
    \endgroup}
\else
  \newcommand{\paperfigure}[1]{\input{#1}}
\fi
\newtheorem{theorem}{Theorem}
\newtheorem{lemma}[theorem]{Lemma}
\newtheorem{proposition}[theorem]{Proposition}
\newtheorem{corollary}[theorem]{Corollary}
\theoremstyle{remark}\newtheorem{remark}[theorem]{Remark}
\newcommand{\enc}{\mathscr E}\newcommand{\R}{\mathbb R}\newcommand{\C}{\mathbb C}
\newcommand{\sq}{\mathrm{sq}}\newcommand{\T}{\mathsf T}\newcommand{\suc}{\succeq}
\hypersetup{pdftitle={Optimal Continuous Analog Encoders for Moment-Based Computation},pdfauthor={Saeed Razavikia}}
\title{Optimal Continuous Analog Encoders for\\Moment-Based Computation}
\author{Saeed Razavikia,~\IEEEmembership{Member,~IEEE} \thanks{Saeed Razavikia is with the Department of Network and Systems Engineering, KTH Royal Institute of Technology, Stockholm, Sweden (e-mail: \href{mailto:sraz@kth.se}{sraz@kth.se}).}}
\begin{document}\maketitle
\begin{abstract}
We study the computation of moment-based functions, including variance, from simultaneous analog transmissions over one complex multiple-access channel use. Each transmitter applies the same continuous encoder to a bounded scalar input and obeys a pointwise peak-power constraint. For at least four transmitters and targets satisfying the stated regularity and strict convexity or concavity assumptions, we prove that exact computation forces the encoder to preserve two specific features: the input and its prescribed nonlinear transformation. This reduces the design of an arbitrary continuous encoder to the geometry of a planar curve. We derive the exact worst-case separation score, formulate globally optimal encoder design as a convex program, and obtain explicit quadratic and quartic curves. A certified numerical procedure handles other smooth features, with an exponential example achieving a global gap below $0.006\%$. A simple projection decoder attains the smallest possible uniform noise-amplification constant. For this decoder and a specified input distribution, we characterize the best small-noise mean-squared error subject to a required separation score. The resulting tradeoff is continuous and convex and admits an explicit scalar characterization for quadratic features. For independent uniform inputs, retaining $90\%$ of optimal separation reduces the small-noise error coefficient by $17.68\%$ and $29.07\%$ with four and sixteen transmitters, respectively. Numerical experiments and deterministic mismatch bounds assess the designs beyond the ideal noiseless model.
\end{abstract}
\begin{IEEEkeywords}Over-the-air computation, analog encoding, multiple-access channels, moment representations, convex optimization, Lipschitz decoding.\end{IEEEkeywords}
\begin{figure*}[t]
 \centering\paperfigure{figures/tikz/system_model}
 \caption{One-use analog computation over a complex multiple-access channel. Every transmitter applies the same continuous encoder to its calibrated input and satisfies the pointwise peak constraint. The channel superposes the symbols; the receiver estimates the prescribed moment-and-mean target from one noisy aggregate. Effective channel gains are unity after compensation, and the payload count excludes calibration and synchronization overhead.}
 \label{fig:system}
\end{figure*}
\section{Introduction}
\IEEEPARstart{F}{unction} computation over multiple-access channels uses the superposition of simultaneous transmissions to evaluate a desired function of distributed observations at a common receiver~\cite{nazer2007,sahin2023survey}. In the conventional analog formulation of over-the-air computation, the observations and the desired function are real-valued. Each transmitter preprocesses its observation, and the receiver applies a postprocessing function to the sum; analog amplitude modulation provides a direct implementation of this principle~\cite{goldenbaum2015,channelcomp2024}. For example, transmitting scaled observations allows the receiver to compute their arithmetic mean. The complex baseband representation of the wireless channel specifies two real signal coordinates and does not require a complex-valued desired function.

Our starting point is function-dependent digital modulation, as developed in ChannelComp~\cite{channelcomp2024}. A digital encoder assigns a constellation point to each quantized observation. Because the channel adds the transmitted symbols, their placement determines which input configurations the receiver can distinguish. For computation, the relevant distances are between received sums associated with different desired function values. Two such sums can be close even when the corresponding function values differ substantially, allowing a small disturbance to cause a large computation error. We therefore seek an input-to-symbol mapping that enables the desired computation and controls its sensitivity to noise under limited transmit power.

We study this design question for real-valued observations through a shared continuous map $\enc:I\to\C$, where $I=[-1,1]\subset\R$. As the observation varies, its transmitted symbol traces a curve in the in-phase and quadrature plane; restricting the map to a finite input alphabet gives a digital constellation. The two coordinates provide two additive signal components within one complex channel use. For variance, a symbol proportional to $s+i s^2$ produces a received sum proportional to $\sum_k s_k+i\sum_k s_k^2$. The receiver can therefore recover the first moment $m$ and second raw moment $q$ and compute $q-m^2$. This example explains the role of the complex-plane curve: it specifies how to convey the quantities needed for a real-valued computation through the two available signal coordinates.

A direct transmission of $s$ and a prescribed preprocessing function $\phi(s)$ establishes achievability, but leaves the optimal encoder geometry undetermined. Increasing the scale of one transmitted quantity can reduce its sensitivity to noise, while the pointwise power constraint limits the scale available to both quantities. We must determine whether exact computation forces every continuous encoder to preserve these two sums, and then optimize their scaling, relative orientation, and common offset. The continuous formulation lets us address both questions over all shared continuous encoders on $I$. The input parameterization remains part of the design: two encoders with the same curve image can produce different received sums for the same observations.

We address these questions for real-valued desired functions $f_{\phi,\psi}:I^K\to\R$ of the form
$$
 f_{\phi,\psi}(\mathbf s)
 =\frac1K\sum_{k=1}^{K}\phi(s_k)
 -\psi\!\left(\frac1K\sum_{k=1}^{K}s_k\right),
 \qquad \mathbf s\in I^K,
$$
where $\phi,\psi:I\to\R$ are prescribed. This class includes variance when $\phi(s)=\psi(s)=s^2$. Section~\ref{sec:model} gives further examples and explains its relation to scalar nomographic representations. We use one shared continuous encoder and one complex channel use. Our design objective is the infimum, over input pairs with different desired function values, of the squared distance between their received sums divided by their squared function difference. The inverse square root of this \emph{squared-separation score} is the smallest uniform disturbance-amplification constant achievable by an exact decoder.

\subsection{Contributions}
We connect exact computation, optimal encoder geometry, and receiver error through the following results.
\begin{enumerate}
\item \emph{Necessary encoder structure.} For $K\ge4$ under the rigidity assumptions in Table~\ref{tab:assumptions}, we prove that every continuous encoder computing $f_{\phi,\psi}$ exactly has the form $\enc(s)=c+us+v\phi(s)$, where $c,u,v\in\C$ and $u,v$ are linearly independent over $\R$. Thus every exact encoder preserves the two sums $S=\sum_k s_k$ and $T=\sum_k\phi(s_k)$. We derive this conclusion from input triples with equal received sums without assuming that the encoder is differentiable. The result permits preprocessing functions $\phi$ with affine subintervals.
\item \emph{Optimal encoding and sharp disturbance bounds.} Under the stronger assumptions in Table~\ref{tab:assumptions}, we express the exact squared-separation score through two endpoint gradients and obtain convex formulations of the peak-constrained encoder design. We derive an equivalent enclosing-ellipse representation of the curve $(s,\phi(s))\in\R^2$. A unified theorem gives explicit optimal encoders for quadratic and quartic preprocessing. We extend the encoder characterization and convex design to functions of two prescribed additive statistics under a common-completion identifiability condition. Polynomial moment interactions, including $mq$, and a nonpolynomial transformation of variance illustrate the extension. For an exponential preprocessing function, we certify a numerical design with a relative global gap below $0.006\%$. A decoder using a fixed linear transformation, scalar clipping, and function evaluation attains the smallest uniform disturbance-amplification constant, including when $u$ and $v$ are nonorthogonal.
\item \emph{A global robustness--accuracy tradeoff.} For a specified input distribution and the constructed decoder, we minimize the leading mean-squared error (MSE) coefficient under circular complex Gaussian noise over all continuous exact encoders meeting a required separation score. A five-variable convex program attains the optimum and determines a unique $2\times2$ Gram matrix of the encoder coefficients $u,v$. The optimal value is continuous and convex in the required score, with an initial plateau followed by a strictly increasing branch. For $\phi(s)=s^2$, we obtain an explicit scalar-root characterization. These results quantify how the optimal curve changes when the designer varies the required uniform robustness and optimizes average small-noise accuracy.
\end{enumerate}
Paired finite-noise experiments, deterministic implementation-error bounds, and a qualitative approximate-rigidity result support the analysis. We provide the detailed optimality proofs and comparisons with other pairwise design criteria in the appendices.

\subsection{Related Work}
Nazer and Gastpar established an information-theoretic framework for computation over multiple-access channels~\cite{nazer2007}. Their compute-and-forward construction uses nested lattice codes to recover linear combinations of transmitted messages~\cite{nazer2011}. For nonlinear real-valued functions, nomographic computation combines local preprocessing, channel summation, and scalar postprocessing. Goldenbaum, Boche, and Sta\'nczak combine this representation with nested lattice coding and analyze computation rates under a block power constraint~\cite{goldenbaum2015}. Our desired function depends on two prescribed real sums, which we encode jointly in the two coordinates of a single complex symbol. We characterize the necessary continuous encoder structure and optimize its sensitivity to disturbances under a pointwise peak constraint.

ChannelComp formulates function-dependent modulation through separation constraints on received superpositions~\cite{channelcomp2024}. Extensions and variants address multiple-antenna computation, blind federated learning, structured digital coding, and constellation sampling~\cite{veccomp2026,blindfl2024,sumcomp2025,pyramidsampling2025}. The original finite-input squared-distance formulation motivates our objective, but constrains the norm of a finite modulation vector rather than each symbol's energy. Noise-aware ChannelComp uses criteria adapted to the noise distribution~\cite{noiseaware2025}; its Gaussian criterion has an uncentered exponential numerator. Analytical optimization also exists for prescribed sum-computing grid constellations with ML or MAP decoding~\cite{towards2025}. Our continuous curve formulation complements these digital designs by establishing an encoder converse and a global squared-separation optimum for the stated classes of nonlinear real-valued functions.

Shared sum representations also underlie Deep Sets~\cite{zaheer2017}. Wagstaff \emph{et al.}\ establish latent-dimension requirements for universal continuous representation and approximation~\cite{wagstaff19,wagstaff22}. Their exact-max argument combines the analysis of equal encoded sums with topology; in particular, \cite[Lemma 17]{wagstaff22} treats a fixed cardinality. These results allow particular desired functions to admit smaller representations. Chen, Chen, and Lu establish representation through multisymmetric polynomial generators and examine singularities of the postprocessing function~\cite{chen2024}.

The closest affine-uniqueness precedent is the restricted Cauchy--Pexider theorem of Charalambides~\cite[Theorem 1.6]{charalambides2013}. For smooth curves without affine subarcs, it identifies affine summands from an additive factorization; its proof regularizes the summands before differentiation. We begin with exact computation of a scalar function and derive the required local factorization from the implication that equal received sums give equal function values. Our assumptions also permit $C^1$ curves with affine subintervals. We retain a self-contained proof of this reduction and do not claim affine uniqueness as a new general principle. Scalar monotone-generator uniqueness for quasi-arithmetic means provides a further classical precedent~\cite[Theorems A and D]{matkowski2015}.

The small-noise objective has the weighted inverse-trace form familiar from A-optimal experimental design. Conic formulations for such criteria and additional resource constraints are classical~\cite{sagnol2010}; semidefinite methods also cover continuous polynomial and rational design spaces~\cite{papp2011}. We apply these tools after identifying the encoder family forced by exact computation. Our feasible matrices arise from fitting the prescribed curve $(s,\phi(s))$ into a peak-power disk through an affine transformation. The resulting robustness--accuracy tradeoff is global for this computation class and the stated decoder and small-noise model. Matrix optimization and enclosing-ellipsoid principles themselves are established~\cite{boyd2004}.

We use classical invariance of domain~\cite{hatcher} and scalar Lipschitz extension~\cite{mcshane1934} to establish the encoder characterization and its operational interpretation.

Section~\ref{sec:model} introduces the channel, desired function class, and objective. Section~\ref{sec:design} combines the necessary encoder structure with optimal peak-constrained design. Section~\ref{sec:examples} presents the quadratic and quartic cases and the extension to two prescribed additive statistics; Section~\ref{sec:decoding} constructs the associated receivers. Section~\ref{sec:pareto} develops the small-noise tradeoff. Sections~\ref{sec:certified} and~\ref{sec:evidence} address certified numerical design and performance. The appendices contain the detailed proofs, additional consequences, and Lean-verification scope.

\section{System Model}\label{sec:model}
\subsection{Multiple-Access Channel}
We consider the sensing architecture in Fig.~\ref{fig:system}. Let $K\in\mathbb N$ denote the number of transmitters and let $I=[-1,1]\subset\R$. Transmitter $k\in\{1,\ldots,K\}$ holds a calibrated real observation $s_k\in I$. We write $\mathbf s=(s_1,\ldots,s_K)^\T\in I^K\subset\R^K$. All transmitters apply the same continuous encoder $\enc:I\to\C$ and transmit simultaneously over one complex channel use. We identify $\C$ with $\R^2$: the real and imaginary components are the in-phase and quadrature coordinates of the transmitted symbol. The inputs and the desired function remain real valued. Unless stated otherwise, our converse and optimal-design results assume $K\ge4$.

For a given $P>0$, we impose the per-transmitter pointwise power constraint
\begin{equation}
 x_k=\enc(s_k),\qquad \sup_{s\in I}|\enc(s)|^2\le P.
 \label{eq:peak-model}
\end{equation}
Thus every admissible observation produces a symbol satisfying the transmit-power limit. We assume synchronized transmissions and unit effective channel gains after compensation. The receiver observes
\begin{equation}
 Y=\sum_{k=1}^{K}x_k+N=r_\enc(\mathbf s)+N,
 \qquad \widehat f(Y)=D(Y),
 \label{eq:mac-model}
\end{equation}
where we define the noiseless superposition $r_\enc:I^K\to\C$ by $r_\enc(\mathbf s)=\sum_{k=1}^K\enc(s_k)$ and use a receiver $D:\C\to\R$. The disturbance $N\in\C$ is arbitrary in our uniform error analysis. For the small-noise analysis, we specify $N\sim\mathcal{CN}(0,N_0)$, independent of $\mathbf s$, with $\mathbb E|N|^2=N_0$. The receiver estimates the desired function directly from the superposition. Calibration and synchronization are assumed resources; their overhead is excluded from the one-use payload count. Section~\ref{sec:implementation} quantifies bounded residual implementation and channel errors.

\subsection{Desired Functions and Examples}
We fix the input calibration and desired function before designing the encoder. For continuous real-valued functions $\phi,\psi:I\to\R$, we consider $f_{\phi,\psi}:I^K\to\R$ defined by
\begin{equation}
 f_{\phi,\psi}(\mathbf s)
 =\frac1K\sum_{k=1}^K\phi(s_k)
   -\psi\!\left(\frac1K\sum_{k=1}^Ks_k\right),
 \qquad \mathbf s\in I^K.
 \label{eq:target}
\end{equation}
We use $\phi$ as a local preprocessing function and $\psi$ as a function of the empirical mean. Define the real sums $S=\sum_k s_k$ and $T=\sum_k\phi(s_k)$, and their normalized statistics $m=S/K\in I$ and $t=T/K\in\R$. Thus $f_{\phi,\psi}(\mathbf s)=t-\psi(m)$. For integer $p\ge1$, let $t_p=K^{-1}\sum_k s_k^p$ be the empirical $p$th raw moment. In particular, $Q=\sum_k s_k^2$ and $q=Q/K=t_2$.

\emph{Mean.} Choosing $\phi(s)=s$ and $\psi\equiv0$ gives $f(\mathbf s)=m$. The encoder $\enc(s)=\sqrt P\,s$ and noiseless receiver $D(y)=\operatorname{Re}(y)/(K\sqrt P)$ compute the mean using one real additive statistic. This example belongs to \eqref{eq:target}, but does not satisfy the nonaffineness and strict-convexity assumptions of our converse. We do not claim that the mean requires two transmitted statistics.

\emph{Variance.} Choosing $\phi(s)=\psi(s)=s^2$ gives the empirical variance
$$
 V(\mathbf s)=q-m^2=\frac1K\sum_{k=1}^K(s_k-m)^2.
$$
For example, $\enc(s)=\sqrt{P/2}\,(s+i s^2)$ obeys \eqref{eq:peak-model}. In the absence of noise, $Y_0=r_\enc(\mathbf s)$ yields
$$
 m=\frac{\operatorname{Re}Y_0}{K\sqrt{P/2}},\qquad
 q=\frac{\operatorname{Im}Y_0}{K\sqrt{P/2}},\qquad V=q-m^2.
$$
This construction explains how two real moments can be conveyed in one complex channel use. Section~\ref{sec:examples} optimizes their scaling and the translation of the transmitted curve.

\emph{Jensen gaps.} If $\psi=\phi$ is convex, \eqref{eq:target} is the empirical Jensen gap. The choices $\phi(s)=s^4$ and $\phi(s)=e^s$ give $t_4-m^4$ and $K^{-1}\sum_k e^{s_k}-e^m$, respectively. The quantity $t_4-m^4$ is not the fourth central moment, which also involves the second and third raw moments. More generally, the boundedness of $\phi$ makes $\enc(s)=\alpha[s+i\phi(s)]$ feasible for $\alpha=\sqrt{P/(1+\|\phi\|_\infty^2)}>0$, so both $S$ and $T$ are sufficient for exact computation. Our converse asks when every exact shared continuous encoder must preserve both sums.

\emph{Relation to nomographic functions.} A scalar nomographic representation has the form
$$
 f(\mathbf s)=\Phi\!\left(\sum_{k=1}^K h_k(s_k)\right),
$$
where $h_k:I\to\R$ are local preprocessing functions and $\Phi$ is a scalar postprocessing function defined on the attainable sums~\cite{goldenbaum2015}. In contrast, our construction uses two real additive statistics, $(S,T)\in\R^2$, followed by $F(m,t)=t-\psi(m)$ with $(m,t)=(S/K,T/K)$. It is therefore a two-coordinate version of nomographic computation, implemented through the two coordinates of one complex symbol. These classes overlap: the mean is already scalar nomographic. Under the assumptions of Theorem~\ref{thm:rigidity}, however, no shared continuous scalar encoder $h:I\to\R$ can compute \eqref{eq:target} exactly, because the theorem forces two linearly independent real coefficient vectors. This conclusion concerns the shared continuous-encoder model; it does not exclude representations with different preprocessing functions at different transmitters or discontinuous preprocessing.

\subsection{Exact Computation and Separation}
For a specified continuous desired function $f:I^K\to\R$, we call $\enc$ \emph{exact} if
\begin{equation}
 \forall\,\mathbf s,\mathbf t\in I^K,\qquad
 r_\enc(\mathbf s)=r_\enc(\mathbf t)
 \ \Longrightarrow\ f(\mathbf s)=f(\mathbf t).
 \label{eq:exactness}
\end{equation}
Equivalently, there is a well-defined noiseless receiver $D_0:\mathcal R_\enc\to\R$, where $\mathcal R_\enc=r_\enc(I^K)\subset\C$, such that $D_0(r_\enc(\mathbf s))=f(\mathbf s)$ for every $\mathbf s\in I^K$. This receiver is continuous: the continuous surjection from the compact set $I^K$ onto the Hausdorff set $\mathcal R_\enc$ is a quotient map. Exactness specifies which function values the superposition must distinguish; the distances between these superpositions determine noise sensitivity.

For a nonconstant $f:I^K\to\R$, we define the squared-separation score
\begin{equation}
 \Lambda_{\sq}(\enc;f)
 =\inf_{\substack{\mathbf s,\mathbf t\in I^K\\f(\mathbf s)\ne f(\mathbf t)}}
 \frac{|r_\enc(\mathbf s)-r_\enc(\mathbf t)|^2}
 {|f(\mathbf s)-f(\mathbf t)|^2}\in[0,\infty).
 \label{eq:score}
\end{equation}
A larger score means that unequal values of the desired function produce better-separated received sums. In particular, $\Lambda_{\sq}>0$ guarantees $|f(\mathbf s)-f(\mathbf t)|\le |r_\enc(\mathbf s)-r_\enc(\mathbf t)|/\sqrt{\Lambda_{\sq}}$ for every $\mathbf s,\mathbf t\in I^K$ and implies exactness. Proposition~\ref{prop:amplification} identifies $1/\sqrt{\Lambda_{\sq}}$ as the smallest uniform noise-amplification constant achievable by a receiver.

We maximize \eqref{eq:score} over the class
$$
 \mathcal E_P=\{\enc\in C(I;\C):\sup_{s\in I}|\enc(s)|^2\le P\},\qquad
 \Lambda_{\sq}^\star(f)=\sup_{\enc\in\mathcal E_P}\Lambda_{\sq}(\enc;f).
$$
Under the assumptions below, Section~\ref{sec:design} proves that every exact encoder has the form $\enc(s)=c+us+v\phi(s)$, giving $r_\enc(\mathbf s)=Kc+uS+vT$. The offset $c$ places the curve inside the peak-power disk; the lengths and relative angle of $u,v$ control the sensitivity of moment recovery. This result reduces the optimization over continuous functions to a finite-dimensional choice of encoder coefficients.

\subsection{Assumptions}
Table~\ref{tab:assumptions} states which assumptions each result requires. We separate the characterization of exact encoders from the stronger conditions needed for a positive sharp separation score. Differentiability only in the interior of $I$ does not ensure finite sensitivity at its endpoints. Strict convexity or concavity of the local preprocessing function $\phi$ suffices for the endpoint rank argument used in the score formula; we do not claim that it is necessary.

Throughout, $C^1(I)$ denotes real-valued continuously differentiable functions with finite continuous one-sided derivatives at the endpoints. We state separately when only interior differentiability is needed.

\begin{table*}[t]
\centering
\setlength{\tabcolsep}{5pt}
\caption{Assumptions and scope of the main results. Throughout, $I=[-1,1]\subset\R$, $K\ge4$, $P>0$, $f:I^K\to\R$, and the shared encoder $\enc:I\to\R^2\cong\C$ is continuous. Optimization results impose the pointwise power constraint \eqref{eq:peak-model}.}
\label{tab:assumptions}
\footnotesize
\begin{minipage}{0.98\textwidth}
\textbf{R (exact-encoder characterization):} $\phi,\psi\in C(I)\cap C^1((-1,1))$; $\phi$ is nonaffine; $\psi$ is strictly convex or strictly concave.\\[2pt]
\textbf{S (sharp separation and design):} R, with $\phi$ strictly convex or strictly concave and $\psi\in C^1(I)$.
\end{minipage}
\vspace{5pt}
\renewcommand{\arraystretch}{1.2}
\begin{tabular}{@{}>{\raggedright\arraybackslash}p{0.27\textwidth}>{\centering\arraybackslash}p{0.07\textwidth}>{\raggedright\arraybackslash}p{0.60\textwidth}@{}}
\toprule
Result & Regime & Scope and additional requirements\\
\midrule
Exact encoders, Theorem~\ref{thm:rigidity} & R & $f=f_{\phi,\psi}$; exactness as in \eqref{eq:exactness}.\\
Optimal separation, Theorem~\ref{thm:sdp} & S & $f=f_{\phi,\psi}$; global optimization over $\mathcal E_P$.\\
Explicit curves, Theorem~\ref{thm:parabola} & S & $\phi(s)=s^2$ or $s^4$; $\psi$ is even.\\
Two-statistic extension, Theorem~\ref{thm:extension} & --- & $\gamma\in C(I;\R^2)\cap C^1((-1,1);\R^2)$; $1,\varphi_1,\varphi_2$ linearly independent; $F\in C(\mathcal M;\R)$; condition~\eqref{eq:completion-id}. Part 2 also requires nonconstant Lipschitz $F$. Formula~\eqref{eq:polyopt} further assumes quadratic $\gamma$, $C^1$ postprocessing near $\mathcal C_2$, and \eqref{eq:general-endpoints}.\\
Projection receiver, Theorem~\ref{thm:decoder} & S & $f=f_{\phi,\psi}$; $\enc(s)=c+us+v\phi(s)$ with $\det[\,u\ v\,]\ne0$; specified strip-projection receiver.\\
Small-noise frontier, Theorem~\ref{thm:pareto} & S & Fixed joint input distribution $\pi$; specified strip decoder; $|m|<1$ almost surely; $Q_\pi\succ0$; independent circular Gaussian noise; any output clipping is locally inactive almost surely.\\
Certified design, Theorem~\ref{thm:certifiedgrid} & S & $\phi\in C^2(I)$; grids contain both endpoints and three fixed points with linearly independent $z(s)$; numerical implementation uses verified function, derivative, and matrix bounds.\\
\bottomrule
\end{tabular}
\end{table*}

\paragraph{Endpoint slopes}
Consider $\phi(s)=s^2$ and $\psi(s)=-\sqrt{1-s}$. Both satisfy R, but $\psi'$ diverges at $1$. Theorem~\ref{thm:rigidity} makes every exact continuous encoder an affine quadratic curve. Compare the two equal-input configurations $s_k=1$ and $s_k=1-h$ for all $k$, where $h\downarrow0$. Their received difference is
$$
\Delta r=K[-hu+(-2h+h^2)v]=O(h),
$$
whereas the difference between their desired function values is $\sqrt h-2h+h^2$. Their squared-separation quotient therefore tends to zero. Every exact encoder has zero score, and every nonexact encoder has zero score by \eqref{eq:score}. Thus R alone permits a zero global optimum and no receiver with finite uniform disturbance amplification.

We therefore first establish which sums exactness requires the transmitters to preserve, and then impose S to optimize robustness.

\section{Encoder Design}\label{sec:design}
We first identify the encoder family required for exact computation and then optimize its geometry under the peak constraint.

\subsection{Exact Encoders}\label{sec:rigidity}
A real-valued desired function $f:I^K\to\R$ might appear to require less information than its two defining statistics $S$ and $T$. The next theorem rules out that possibility under the rigidity assumptions: computing $f$ exactly for every input configuration forces the shared encoder to preserve both sums. This gives the converse needed to compare a proposed curve with all continuous alternatives.

\begin{theorem}
\label{thm:rigidity}
Let $K\ge4$. Assume $\phi\in C(I)\cap C^1((-1,1))$ is nonaffine and $\psi\in C(I)\cap C^1((-1,1))$ is strictly convex or strictly concave. A continuous encoder $\enc:I\to\R^2\cong\C$ exactly computes $f_{\phi,\psi}:I^K\to\R$ in \eqref{eq:target}, in the sense of \eqref{eq:exactness}, if and only if
\begin{equation}
 \enc(s)=c+u s+v\phi(s),
 \qquad s\in I,\quad c,u,v\in\R^2,
 \label{eq:affine}
\end{equation}
with $u,v$ linearly independent, equivalently $\det[\,u\ v\,]\ne0$.
\end{theorem}
The proof appears in Appendix~\ref{app:rigidity-proof}. Its key step uses a fourth transmitter to show that equal encoded triple sums must preserve both sums $S$ and $T$. The receiver can then recover $S,T$ from $r_\enc-Kc=uS+vT$ and evaluate $f_{\phi,\psi}$. Thus exact computation leaves the scale, orientation, and offset of the preprocessing curve $(s,\phi(s))$ available for noise-aware design.

The classification uses one complex channel use and $K\ge4$. It does not establish the necessity of \eqref{eq:affine} for unrestricted continuous encoders when $K=3$.\footnote{For $K=3$, a collision between two encoded triples implies only $T-\widetilde T=3[\psi(S/3)-\psi(\widetilde S/3)]$; changes in the two statistics can cancel in this identity. For $K\ge4$, we append the same fourth input $w\in I$ and $K-4$ zeros to both triples. Exactness then gives $T-\widetilde T=K[\psi((S+w)/K)-\psi((\widetilde S+w)/K)]$ for every $w\in I$. Comparing $w=0$ and $w=1/2$, strict convexity or concavity forces $S=\widetilde S$, and hence $T=\widetilde T$; see Lemma~\ref{lem:collision}. There is no fourth input to vary when $K=3$, so this proof does not establish preservation of both sums in that case. The affine form with independent $u,v$ remains sufficient for exact computation when $K=3$; only its necessity is unresolved here.} Its assumptions permit unbounded endpoint derivatives and therefore do not alone imply positive separation.

\subsection{Optimal Design}\label{sec:synthesis}
We now assume $\phi\in C(I)\cap C^1((-1,1))$ and $\psi\in C^1(I)$, each strictly convex or strictly concave. Finite endpoint slopes of the mean-dependent term provide a uniform disturbance guarantee; $\phi$ needs differentiability only in the interval interior. For the encoder coefficients in \eqref{eq:affine}, define
\begin{align*}
 G&=\begin{pmatrix}A&C\\ C&B\end{pmatrix},\quad
 w_\pm=\begin{pmatrix}-\psi'(\pm1)\\1\end{pmatrix},\\
 A&=|u|^2,\quad B=|v|^2,\quad C=\operatorname{Re}(u\overline v).
\end{align*}
Here $A,B$ are the squared lengths of the two feature directions and $C$ is their real inner product. The Gram matrix $G$ therefore records the part of the transmit geometry that affects received distances; $C=0$ means orthogonal directions. The vectors $w_\pm$ describe the target sensitivity at the extreme input means.

\begin{theorem}
\label{thm:sdp}
Let $K\ge4$ and $P>0$, and impose the assumptions of this subsection. For every encoder \eqref{eq:affine} with independent $u,v$,
\begin{equation}
 \Lambda_{\sq}(\enc;f_{\phi,\psi})
 =\frac{K^2}{\max\{w_-^\T G^{-1}w_-,\;w_+^\T G^{-1}w_+\}}.
 \label{eq:exactscore}
\end{equation}
The infimum is approached at arbitrarily small received distance. The global optimum over all continuous peak-admissible encoders is $K^2\lambda_\star$, where $\lambda_\star$ is the optimum of
\begin{equation}
\begin{aligned}
 \text{maximize}\quad &\lambda\\
 \text{over}\quad &H\in\mathbb S^3,\quad\lambda\ge0,\\
 \text{subject to}\quad
 &H\suc0,\\
 &z(s)^\T H z(s)\le P\quad\text{for every }s\in[-1,1],\\
 &G-\lambda w_-w_-^\T\suc0,\\
 &G-\lambda w_+w_+^\T\suc0.
\end{aligned}
\label{eq:sdp}
\end{equation}
where $z(s)=(1,s,\phi(s))^\T$ and $G$ is the lower-right $2\times2$ block of $H$. This is a convex semidefinite problem with seven real variables and a continuum of scalar peak constraints. Its optimum is attained and yields an exact planar encoder.
\end{theorem}
See Appendix~\ref{app:design-proof} for the score calculation and the proof of optimality and attainment. Equation~\eqref{eq:exactscore} separates two design roles: $\phi$ determines the feasible peak geometry, while $\psi$ enters through its two extreme slopes.

To implement a solution, write $H=\left(\begin{smallmatrix}h&\ell^\T\\\ell&G\end{smallmatrix}\right)$, choose $L^{\T}L=G$, and set $c=(L^\T)^{-1}\ell$. The columns of $L$ give $u,v$. Reducing $h$ to $\ell^{\T}G^{-1}\ell$ preserves the score and can only lower the peak energy, so a planar encoder always suffices.

The peak constraint must hold between samples. Section~\ref{sec:examples} gives explicit curves with analytic certificates, and Section~\ref{sec:certified} provides certified numerical design for other smooth features.

\subsection{Enclosing Ellipses}
\label{sec:inverse-synthesis}
Inverse-metric coordinates give a smaller equivalent formulation and make the design geometry explicit. Write $g(s)=(s,\phi(s))^\T$. Let $\mathcal C_P$ be the set of pairs $X\in\mathbb S^2$, $a\in\R^2$ satisfying
\begin{equation}
 \begin{pmatrix}P&(g(s)-a)^\T\\g(s)-a&X\end{pmatrix}
 \suc0\quad(s\in I).
 \label{eq:ellipse-class}
\end{equation}
Every member has $X\succ0$. Indeed, a null vector of $X$ would be orthogonal to $g(s)-a$ for every $s$, confining the nonaffine feature graph to an affine line. The Schur complement shows that $\mathcal C_P$ describes ellipses
$$
 (g(s)-a)^\T X^{-1}(g(s)-a)\le P
$$
containing the feature curve. We reconstruct an encoder by choosing $\mathsf U^\T\mathsf U=X^{-1}$ and setting $\enc(s)=\mathsf U(g(s)-a)$. Conversely, every exact encoder defines $X=G^{-1}$ and $a=-\mathsf U^{-1}c$.

We may equivalently solve the six-variable convex program
\begin{equation}
\begin{aligned}
 t_\star=\min_{X,a,t}\quad &t\\
 \text{subject to}\quad &(X,a)\in\mathcal C_P,\\
 &w_-^\T Xw_-\le t,\qquad w_+^\T Xw_+\le t.
\end{aligned}
 \label{eq:inverse-synthesis}
\end{equation}
and recover the optimal score as $K^2/t_\star$. Corollary~\ref{cor:inverse-synthesis} in Appendix~\ref{app:design-proof} establishes equivalence and attainment.

This is an enclosing-ellipsoid formulation of the same computation problem. We minimize target-sensitive widths rather than volume, using classical convex geometric tools~\cite[Chap. 8]{boyd2004}. The Gram form remains convenient for the numerical certificate in Section~\ref{sec:certified}; the inverse form makes the prior-dependent objective in Section~\ref{sec:pareto} linear.

\section{Optimal Curves}\label{sec:examples}

The convex design problem gives explicit solutions for several useful desired functions. We first present quadratic and quartic preprocessing functions $\phi(s)$ in \eqref{eq:target}, including variance and the fourth-order Jensen gap. We then extend the framework to general functions of two prescribed additive statistics. Figure~\ref{fig:curves} illustrates the resulting modulation curves and the source values at which they use the full allowed peak power.

\subsection{Quadratic and Quartic Preprocessing Functions}
\label{sec:quadratic-quartic}
For $p\in\{2,4\}$, write $\phi_p(s)=s^p$. Let $\psi:I\to\R$ be even, continuously differentiable on $I=[-1,1]$, and strictly convex or strictly concave. Its endpoint slope $M=|\psi'(1)|$ is positive. For $\phi(s)=s^2$, define
$$
 b=\frac{2\bigl(\sqrt{M^2+8}+M\bigr)}{\sqrt{M^2+8}+3M},
 \qquad a=\sqrt{b(2-b)}.
$$
Here $1<b<2$. For $\phi(s)=s^4$, let $\rho\in(0,1/3)$ be the unique solution of
\begin{equation}
 8\rho^2(1+\rho)(1+3\rho)=M^2(1-3\rho),
 \label{eq:quarticroot}
\end{equation}
and define
\begin{equation}
\begin{aligned}
 d_\rho&=1+2\rho+9\rho^2,\qquad
 a_\rho=\sqrt{\frac{8\rho}{d_\rho}},\\
 b_\rho&=\frac{2}{(1+\rho)\sqrt{d_\rho}},\qquad
 c_\rho=-\frac{1+2\rho+3\rho^2}{(1+\rho)\sqrt{d_\rho}}.
\end{aligned}
\label{eq:quarticcoefficients}
\end{equation}
The root is well defined because the left side of \eqref{eq:quarticroot} increases strictly from zero, whereas the right side decreases strictly to zero. The corresponding normalized separation constants are
$$
 \lambda_2=b^2(b-1),\qquad
 \lambda_4=\frac{1}{M^2/a_\rho^2+1/b_\rho^2}
 =\frac{8\rho}{d_\rho\,[M^2+2\rho(1+\rho)^2]}.
$$

\begin{theorem}[Quadratic and quartic encoders]
\label{thm:parabola}\label{thm:quartic}
Let $K\ge4$ be an integer and let $P>0$. Under the preceding assumptions on $\psi$, the desired functions $f_{\phi,\psi}:I^K\to\R$ in \eqref{eq:target} admit the following globally optimal encoders among all continuous maps $\enc:I\to\C$ satisfying $|\enc(s)|^2\le P$ for every $s\in I$:
for $\phi(s)=s^2$,
\begin{equation}
\begin{aligned}
 \enc_2^\star(s)&=\sqrt P\,[a s+i(bs^2-1)],\qquad s\in I,\\
 \Lambda_{\sq}^\star(f_{\phi_2,\psi})&=K^2P\lambda_2;
\end{aligned}
 \label{eq:parabola}
\end{equation}
for $\phi(s)=s^4$,
\begin{equation}
\begin{aligned}
 \enc_4^\star(s)&=\sqrt P\,[a_\rho s+i(b_\rho s^4+c_\rho)],\qquad s\in I,\\
 \Lambda_{\sq}^\star(f_{\phi_4,\psi})&=K^2P\lambda_4.
\end{aligned}
 \label{eq:quarticscore}
\end{equation}
\end{theorem}
Appendix~\ref{app:quartic-proof} proves both claims and verifies the peak constraint on the entire interval. For either preprocessing function, $\psi$ affects the optimal coefficients only through $M$.

For variance, $\phi(s)=\psi(s)=s^2$, so $M=2$, $b=2/\sqrt3$, and
$$
 \Lambda_{\sq}^\star(V)=\frac{4K^2P}{9+6\sqrt3}.
$$
For the fourth-order Jensen gap, $\phi(s)=\psi(s)=s^4$, so $M=4$ and $\lambda_4\simeq0.058533507$. The quadratic curve reaches the peak power at $s=-1,0,1$; the quartic curve reaches it at $s=\pm1$ and $s=\pm\sqrt\rho$. Thus changing the preprocessing function changes both the modulation curve and its active peak constraints.

The theorem also covers $\psi(m)=m^{2r}$ for every integer $r\ge1$, $\psi(m)=-\eta m^{2r}$ for $\eta>0$, and $\psi(m)=m^2+\eta m^4$ for $\eta\ge0$. In the last case, $M=2+4\eta$. Strict convexity allows a vanishing second derivative at isolated points; in particular, it includes $m^4$.

\begin{figure*}[t]
\centering\paperfigure{figures/tikz/optimal_curves}
\caption{Optimal encoders and source values attaining the peak power for $K\ge4$, with $\bar\lambda^\star=\Lambda_\sq^\star/(K^2P)$. The baselines are $\sqrt P[s+i(s^2-1)]$ for variance and the moment product $mq$, and $(2\sqrt P/\sqrt5)[s+i(s^4-1/2)]$ for the quartic Jensen gap. All curves satisfy the same pointwise peak constraint.}
\label{fig:curves}\end{figure*}

\subsection{Extension to Two Prescribed Additive Statistics}
\label{sec:polynomial}\label{sec:extension}
The preceding family combines two statistics in the specific form $t-\psi(m)$. Functions such as the moment product $mq$ also use two sums, but combine them differently. We now give a common characterization and design theorem that includes both forms, with polynomial functions of moments as one example. Let $J=(-1,1)$, and fix
\begin{equation}
\begin{aligned}
 \gamma(s)&=(\varphi_1(s),\varphi_2(s))^\T\in\R^2,\qquad
 \boldsymbol\mu_\gamma(\mathbf s)=\frac1K\sum_{k=1}^K\gamma(s_k),\\
 f(\mathbf s)&=F(\boldsymbol\mu_\gamma(\mathbf s)),\qquad
 \mathcal M=\boldsymbol\mu_\gamma(I^K),\quad F\in C(\mathcal M;\R).
\end{aligned}
 \label{eq:general-function}
\end{equation}
Here $\varphi_1,\varphi_2:I\to\R$ are prescribed local preprocessing functions, and $F$ combines their averages at the receiver. We assume $\gamma\in C(I;\R^2)\cap C^1(J;\R^2)$ and linear independence of the three real functions $1,\varphi_1,\varphi_2$; equivalently, $\gamma(I)$ is not contained in an affine line. The original family uses $\gamma(s)=(s,\phi(s))^\T$ and $F(a,b)=b-\psi(a)$. Other choices include the variance of squared observations, $\gamma(s)=(s^2,s^4)^\T$ and $F(a,b)=b-a^2$.

If we require only continuity and allow $\gamma$ to be chosen, \eqref{eq:general-function} describes precisely the continuous desired functions that admit identifiable computation with a shared continuous planar encoder. Such an encoder supplies its two coordinates as $\gamma$; conversely, a scaled version of any bounded continuous $\gamma$ transmits the two sums. Prescribing $\gamma$ raises a further question: must every encoder permitting identifiable computation of $f$ preserve these particular sums? This requires $F$ to distinguish them sufficiently. For example, $F(a,b)=a$ ignores the second sum.

We state this additional condition using common completions of three inputs. It is a property of the prescribed statistics and $F$, separate from the identifiability condition \eqref{eq:exactness} for an encoder. Define
$$
 T_3(x,y,z)=\gamma(x)+\gamma(y)+\gamma(z),\qquad
 \Omega_3=\{T_3(x,y,z):(x,y,z)\in J^3,\ \operatorname{rank}DT_3(x,y,z)=2\}.
$$
This set is nonempty and open. Fix a filler input $s_\circ\in J$ and put $d_\circ=(K-4)\gamma(s_\circ)$. For $p\in\Omega_3$ and $x\in J$, define the completed function value $R_x(p)=F((p+\gamma(x)+d_\circ)/K)$. Its argument belongs to the actual attainable set $\mathcal M$. We require a dense open subset $\mathcal O\subseteq\Omega_3$ such that each $p_0\in\mathcal O$ has an open neighborhood $B\subseteq\Omega_3$ satisfying
\begin{equation}
 \forall p,\widetilde p\in B,\qquad
 [R_x(p)=R_x(\widetilde p)\ \text{for every }x\in J]
 \ \Longrightarrow\ p=\widetilde p.
 \label{eq:completion-id}
\end{equation}
Thus nearby sums can be distinguished through their desired function values under common remaining inputs. These comparisons are used only in the proof and do not require additional transmissions. For $\gamma(s)=(s,\phi(s))^\T$ and $F(a,b)=b-\psi(a)$, Lemma~\ref{lem:collision} verifies this condition under the original rigidity assumptions.

\begin{theorem}[Extension to two additive statistics]
\label{thm:extension}\label{thm:polynomial}
Let $K\ge4$, $P>0$, and let \eqref{eq:general-function} satisfy the regularity, independence, and common-completion condition stated above.
\begin{enumerate}
\item A continuous encoder $\enc:I\to\R^2\cong\C$ satisfies the identifiability condition \eqref{eq:exactness} for $f$ if and only if
$$
 \enc(s)=c+\mathsf U\gamma(s),\qquad s\in I,\quad
 c\in\R^2,\quad\mathsf U\in\R^{2\times2},\quad\det\mathsf U\ne0.
$$
\item If $F$ is additionally nonconstant and Lipschitz on $\mathcal M$ in the Euclidean norm, then $\Lambda_\sq^\star(f)=K^2\lambda_\gamma^\star$, where $\lambda_\gamma^\star>0$ is the attained optimum of
\begin{equation}
\begin{aligned}
 \text{maximize}\quad &\lambda\\
 \text{over}\quad &H\in\mathbb S^3,\quad\lambda\ge0,\\
 \text{subject to}\quad &H\succeq0,\\
 &z_\gamma(s)^\T H z_\gamma(s)\le P\qquad(s\in I),\\
 &(\mu-\nu)^\T G(\mu-\nu)\ge\lambda|F(\mu)-F(\nu)|^2
       \qquad(\mu,\nu\in\mathcal M).
\end{aligned}
 \label{eq:general-sdp}
\end{equation}
Here $z_\gamma(s)=(1,\gamma(s)^\T)^\T\in\R^3$ and $G$ is the lower-right $2\times2$ block of $H$. Every optimizer yields a planar encoder that satisfies the identifiability condition for $f$ and the peak constraint.
\end{enumerate}
\end{theorem}
Appendix~\ref{app:polynomial} proves the theorem. Program~\eqref{eq:general-sdp} has the same seven scalar variables as \eqref{eq:sdp}; the general postprocessing function requires constraints over attainable pairs in place of the two endpoint inequalities. The finite-$K$ set $\mathcal M$ need not be convex, so replacing it by its convex hull need not preserve the optimum. The theorem characterizes the exact program; finite discretization requires a separate accuracy certificate.

\emph{Verifying the completion condition.} A useful sufficient condition is that $F$ extend to a real analytic function on all of $\R^2$ and satisfy
\begin{equation}
 \forall\ell\in\R^2,\qquad
 [\ell^\T\nabla F\equiv0\text{ on }\R^2]\ \Longrightarrow\ \ell=0.
 \label{eq:analytic-nondegenerate}
\end{equation}
For polynomial $F$, this excludes precisely functions of one linear combination, $F(z)=h(w^\T z)$ with $w\ne0$ and univariate polynomial $h$. Arbitrary coefficient signs, mixed parity, and polynomial degree are allowed for this encoder classification. Appendix~\ref{app:polynomial} proves the sufficient condition. The Lipschitz assumption for design holds, for example, whenever $F$ extends to a $C^1$ function on a neighborhood of $\operatorname{conv}\gamma(I)$.

\emph{An explicit quadratic specialization.} For $\gamma(s)=(s,s^2)^\T$, retain \eqref{eq:completion-id} and suppose $F$ is $C^1$ on a neighborhood of $\mathcal C_2=\{(m,q)\in\R^2:m^2\le q\le1\}$. Set $w_\pm=\nabla F(\pm1,1)$. If constants $\alpha,\beta>0$ satisfy
\begin{equation}
\begin{aligned}
 &|F_m(m,q)|\le\alpha,\quad |F_q(m,q)|\le\beta
       \qquad((m,q)\in\mathcal C_2),\\
 &\{w_+,-w_+,w_-,-w_-\}
   =\{(\epsilon_1\alpha,\epsilon_2\beta)^\T:\epsilon_1,\epsilon_2\in\{-1,1\}\},
\end{aligned}
 \label{eq:general-endpoints}
\end{equation}
then the endpoint argument and quadratic peak-power certificate give an explicit optimum. Define
$$
 b_{\alpha,\beta}
 =\frac{2\bigl(\sqrt{\alpha^2+8\beta^2}+\alpha\bigr)}
 {\sqrt{\alpha^2+8\beta^2}+3\alpha},\qquad
 a_{\alpha,\beta}=\sqrt{b_{\alpha,\beta}(2-b_{\alpha,\beta})}.
$$
An optimal encoder over $\mathcal E_P$ and its score are
\begin{equation}
\begin{aligned}
 \enc^\star(s)&=\sqrt P\,[a_{\alpha,\beta}s+i(b_{\alpha,\beta}s^2-1)],\qquad s\in I,\\
 \Lambda_{\sq}^\star(f)&=\frac{K^2P}{\beta^2}
 b_{\alpha,\beta}^2(b_{\alpha,\beta}-1).
\end{aligned}
 \label{eq:polyopt}
\end{equation}
The proof in Appendix~\ref{app:polynomial} establishes sharpness using attainable perturbations near the two consensus inputs. Conditions~\eqref{eq:general-endpoints} are additional sufficient assumptions for this explicit formula; they are not required by the general design program.

\emph{Polynomial functions of the first two moments.} Polynomial functions of moments provide an example of this specialization:
\begin{equation}
 f(\mathbf s)=F(m,q),\qquad
 F(m,q)=\sum_{(i,j)\in\mathcal I}a_{ij}m^i q^j,
 \label{eq:polyclass}
\end{equation}
where $\mathcal I\subset\mathbb Z_{\ge0}^2$ is finite, $a_{ij}\ge0$, and all exponents $i$ with $a_{ij}>0$ have a common parity $\epsilon\in\{0,1\}$. If
$$
 \alpha=F_m(1,1)=\sum_{(i,j)\in\mathcal I}i a_{ij}>0,\qquad
 \beta=F_q(1,1)=\sum_{(i,j)\in\mathcal I}j a_{ij}>0,
$$
then both \eqref{eq:completion-id} and \eqref{eq:general-endpoints} hold, and \eqref{eq:polyopt} follows. The sign and parity restrictions thus verify the explicit optimum for this example family, while Theorem~\ref{thm:extension} covers more general functions.

For the moment product,
$$
 f(\mathbf s)=mq
 =\frac{1}{K^2}\left(\sum_{k=1}^K s_k\right)
 \left(\sum_{k=1}^K s_k^2\right),
$$
we have $\alpha=\beta=1$, which gives the optimal curve and score:
$$
 \enc^\star(s)=\sqrt P\left[\frac{2\sqrt2}{3}s
 +i\left(\frac43s^2-1\right)\right],\qquad
 \Lambda_{\sq}^\star(f)=\frac{16}{27}K^2P.
$$
The channel provides the same two sums $S$ and $Q$ as in variance computation, while their product requires different optimal scales. Further examples are $F(m,q)=m^2q+\eta q^2+\zeta m^4$ for $\eta,\zeta\ge0$, with $\alpha=2+4\zeta$ and $\beta=1+2\eta$, and $F(m,q)=m^3q^2+mq^4$, with $\alpha=4$ and $\beta=6$.

\emph{A nonpolynomial example.} Let $V=q-m^2$ and $f(\mathbf s)=1-e^{-V}$. The transformation $v\mapsto1-e^{-v}$ is injective on $[0,1]$, so the same encoders permit identifiable computation of this desired function and of variance. Moreover, $\nabla F=e^{-V}(-2m,1)^\T$ satisfies \eqref{eq:general-endpoints} with $\alpha=2$ and $\beta=1$. Therefore its optimal encoder and squared-separation score are also those of variance. This example extends the explicit result beyond polynomial postprocessing without introducing a different transmit curve.

Proposition~\ref{prop:amplification} below applies to the uniform disturbance guarantee for the general extension. The explicit strip receiver and high-SNR MSE frontier retain their stated assumptions on the original family $f_{\phi,\psi}$.

\section{Decoding}\label{sec:decoding}
The preceding designs separate noiseless received points as effectively as the peak constraint permits. To turn this geometric property into a receiver guarantee, we must also define the decoder outside the noiseless received set, where additive noise can move the observation. A polynomial inverse may grow faster than linearly there. We first relate squared separation to uniform noise amplification, then construct a direct decoder that attains the resulting bound.

% \subsection{Uniform Error Bound}
For a fixed encoder, define $\mathcal A(\enc;f)$ as the infimum of constants $L\ge0$ for which one decoder $\widehat f:\R^2\to\R$ satisfies
\begin{equation}
 |\widehat f(r_\enc(\mathbf s)+n)-f(\mathbf s)|\le L|n|
 \quad\forall\mathbf s\in I^K,\ \forall n\in\R^2.
 \label{eq:uniform}
\end{equation}
Set $\mathcal A=+\infty$ if no finite constant exists. This decoder is fixed for every input and disturbance magnitude.

\begin{proposition}\label{prop:amplification}
Let $I=[-1,1]$, $K\in\mathbb N$ with $K\ge1$, $\enc:I\to\R^2$ be continuous, and $f:I^K\to\R$ be continuous and nonconstant. Then
$$\mathcal A(\enc;f)=\Lambda_\sq(\enc;f)^{-1/2},$$
with $0^{-1/2}=+\infty$. A continuous decoder attains every finite value.
\end{proposition}
\begin{proof}
The choice $n=0$ enforces exactness. Using one noiseless received point as the observation under another input gives $|\Delta f|\le L|\Delta r|$, which proves the lower bound. Conversely, a finite supremum of $|\Delta f|/|\Delta r|$ defines a Lipschitz receiver function $d_0:\mathcal R_\enc\to\R$ on the compact noiseless received set $\mathcal R_\enc$. McShane's construction~\cite{mcshane1934},
$$\widehat f(y)=\inf_{z\in\mathcal R_\enc}\{d_0(z)+L|y-z|\},$$
extends $d_0$ with the same constant. The triangle inequality proves global Lipschitz continuity, and exactness gives \eqref{eq:uniform}. Clipping the scalar output to any interval containing $f(I^K)$ preserves the guarantee.
\end{proof}
The extension principle is classical; our encoder optimum identifies the smallest achievable extension constant. Evaluating the infimum over the full received set is unnecessary for implementation. We instead use the recovered statistics $(m,t)$: the sample mean $m$ must lie in $[-1,1]$, and enforcing this constraint gives an explicit projection decoder. Write $\enc(s)=c+us+v\phi(s)$ and let $\mathsf U=[u\ v]\in\R^{2\times2}$ be invertible, with $G=\mathsf U^\T\mathsf U$. Define
$$
 (\widetilde m,\widetilde t)^\T=\frac1K \mathsf U^{-1}(y-Kc).
$$
We project in metric $G$ onto the strip $\mathcal S=[-1,1]\times\R$. Its closed-form solution gives the decoder
\begin{equation}\begin{aligned}
 \widehat m&=\operatorname{clip}_{[-1,1]}(\widetilde m),\\
 \widehat t&=\widetilde t+\frac{G_{12}}{G_{22}}(\widetilde m-\widehat m),\\
 \widehat f(y)&=\widehat t-\psi(\widehat m).
\end{aligned}\label{eq:stripdecoder}\end{equation}
Thus decoding requires a fixed linear transformation, one scalar clipping operation, and evaluation of $\widehat t-\psi(\widehat m)$. Nonorthogonal coefficient vectors $u,v\in\R^2$ require the correction in $\widehat t$.

\begin{theorem}\label{thm:decoder}
Under the assumptions of Theorem~\ref{thm:sdp}, let $\enc(s)=c+us+v\phi(s)$ with $c,u,v\in\R^2$ and linearly independent $u,v$. The receiver function $\widehat f:\R^2\to\R$ in \eqref{eq:stripdecoder} exactly computes $f_{\phi,\psi}:I^K\to\R$ and attains the smallest possible uniform constant for this encoder
\begin{equation}
 L_\enc=\frac1K\sqrt{\max_{\sigma\in\{-,+\}}w_\sigma^\T G^{-1}w_\sigma}
 =\Lambda_\sq(\enc;f)^{-1/2}.
 \label{eq:decoderconstant}
\end{equation}
An additional scalar projection onto an interval containing $f(I^K)$ preserves this conclusion.
\end{theorem}
\begin{proof}
For $x=(m,t)^\T$ and $\widetilde x=(\widetilde m,\widetilde t)^\T$, Euclidean projection onto the received strip $Kc+K\mathsf U\mathcal S$ minimizes $K^2(x-\widetilde x)^\T G(x-\widetilde x)$. For fixed $m$, differentiation in $t$ gives
$$t=\widetilde t-G_{12}(m-\widetilde m)/G_{22}.$$
Substitution leaves the positive Schur-complement coefficient $G_{11}-G_{12}^2/G_{22}$ multiplying $(m-\widetilde m)^2$, so clipping $m$ proves \eqref{eq:stripdecoder}. Euclidean projection onto this closed convex received set is nonexpansive.

The receiver postprocessing function $F:\mathcal S\to\R$, $F(m,t)=t-\psi(m)$, has gradient $(-\psi'(m),1)^\T$. The endpoint-slope argument for \eqref{eq:exactscore} bounds its squared dual norm in metric $G$ by $K^2L_\enc^2$. Integration along segments proves that the corresponding receiver function on $Kc+K\mathsf U\mathcal S\subset\R^2$ is $L_\enc$-Lipschitz. Its composition with nonexpansive projection preserves this constant and fixes all noiseless points. The exact score and Proposition~\ref{prop:amplification} prove optimality.
\end{proof}

\begin{figure}[t]
 \centering\paperfigure{figures/tikz/decoder_geometry}
 \caption{Projection decoding for the optimal variance encoder at $K=4$, with $b=2/\sqrt3$ and $a=\sqrt{b(2-b)}$. The attained received region lies inside the strip $\mathcal S_r=\{(x,z)\in\R^2:|x|\le a\}$. We normalize all received coordinates by $K\sqrt P$ and use equal axis scales. The decoder projects $y$ onto the strip and evaluates $\widehat q-\widehat m^2$; the projected point need not be attainable.}
 \label{fig:decoder-geometry}
\end{figure}
Figure~\ref{fig:decoder-geometry} illustrates why the containing strip is sufficient. For variance at $K=4$, the exact moment region has lower boundary $q=m^2$ and upper boundary $q=[3+(4m-j)^2]/4$, where $j\in\{-3,-1,1,3\}$ and $|4m-j|\le1$. The upper boundary follows by maximizing the convex second moment on a fixed-mean cube slice: a vertex has at most one nonendpoint coordinate. The slice is connected, so every intermediate second moment is attained. This region need not be convex, whereas its containing strip admits the explicit nonexpansive projection.

The strip contains every attainable moment point, so the decoder needs no exact finite-$K$ moment-region characterization. For the local preprocessing function $\phi(s)=s^p$, $p\in\{2,4\}$, a tighter alternative is
$$\mathcal C_p=\{(m,t):-1\le m\le1,\ m^p\le t\le1\}.$$
Jensen's inequality proves containment, and the set is convex. Projection onto its received image has the same guarantee, although its implementation requires convex optimization. It is a containing convex hull, not generally the exact equal-weight finite-$K$ moment set.

For the optimal orthogonal map $\sqrt P[a s+i(bs^p+c_0)]$, compute
\begin{align*}
 \widehat m&=\operatorname{clip}_{[-1,1]}\frac{\operatorname{Re}y}{Ka\sqrt P},\\
 \widehat t&=\frac{\operatorname{Im}y/(K\sqrt P)-c_0}{b}.
\end{align*}
The variance and fourth-order Jensen-gap decoders are $\operatorname{clip}_{[0,1]}(\widehat t-\widehat m^p)$, with the appropriate $p$ and coefficients. Both desired functions take values in $[0,1]$.

\subsection{Moment Interactions and Operational Limits}
The same construction also applies when the receiver combines the recovered moments through a different postprocessing function. For the desired function $f(\mathbf s)=m q$, the containing rectangle $[-1,1]\times[0,1]$ and the optimal orthogonal curve give
\begin{equation}\begin{aligned}
 \widehat m&=\operatorname{clip}_{[-1,1]}\frac{3\operatorname{Re}y}{2\sqrt2K\sqrt P},\\
 \widehat q&=\operatorname{clip}_{[0,1]}\frac34\left(\frac{\operatorname{Im}y}{K\sqrt P}+1\right),\\
 \widehat f_{mq}(y)&=\widehat m\widehat q.
\end{aligned}\label{eq:mqdecoder}\end{equation}
The rectangle's received projection is coordinatewise clipping because the directions are orthogonal. The dual gradient norm of $mq$ is at most $\sqrt{1/a^2+1/b^2}=3\sqrt3/4$ for $a=2\sqrt2/3$, $b=4/3$. Hence this decoder attains $L_{mq}=3\sqrt3/(4K\sqrt P)$. For the restricted polynomial class, the endpoint partial-derivative bounds hold on the same rectangle; projection followed by polynomial postprocessing likewise attains its sharp constant.


These decoders control disturbances of every magnitude with one fixed rule. The raw variance inverse is quadratic in $\operatorname{Re}y$, and the fourth-order inverse contains its fourth power. Neither is globally Lipschitz on the observation plane. Projection is therefore part of the noise guarantee. For a fixed encoder and known disturbance radius $\tau\ge0$, define
$$\omega_\enc(d)=\max_{\substack{\mathbf s,\mathbf t\in I^K\\|r_\enc(\mathbf s)-r_\enc(\mathbf t)|\le d}}|f(\mathbf s)-f(\mathbf t)|,\qquad d\ge0.$$
The minimum worst-case absolute error at radius $\tau$ equals $\omega_\enc(2\tau)/2$. Two configurations separated by at most $2\tau$ can give the same midpoint observation, proving the lower bound. A matching decoder outputs the midpoint of the smallest and largest values of $f(\mathbf s)$ consistent with the observation and radius. Compactness gives these extrema. Their maximum is upper semicontinuous and minimum lower semicontinuous on the compact possible-observation set, so the midpoint rule has a Borel measurable extension. This criterion depends on the full ambiguity profile, rather than one constant for every disturbance size.

For any finite-second-moment noise, \eqref{eq:uniform} implies
$$\mathbb E|\widehat f(Y)-f(\mathbf s)|^2\le L_\enc^2\mathbb E|N|^2,$$
also after averaging over an input prior. Under $N\sim\mathcal{CN}(0,N_0)$, the bound is $L_\enc^2N_0$. This is a bound for a specified decoder, not unrestricted Gaussian-MSE optimality.

\subsection{Implementation Errors}
\label{sec:implementation}
The uniform guarantee also controls encoder approximation and residual channel-compensation errors. Let the implemented map $\widetilde\enc$ satisfy
$$
 \sup_{s\in I}|\widetilde\enc(s)|\le\sqrt P,
 \qquad
 \sup_{s\in I}|\widetilde\enc(s)-\enc(s)|\le\epsilon\sqrt P,
$$
and let the residual channel coefficient of transmitter $k$ be $1+d_k$. The observation is
$$
 Y=\sum_{k=1}^K(1+d_k)\widetilde\enc(s_k)+N.
$$
Applying the ideal decoder and using \eqref{eq:uniform} gives
\begin{equation}
 |\widehat f(Y)-f(\mathbf s)|
 \le L_\enc\left(|N|+K\epsilon\sqrt P
       +\sqrt P\sum_k|d_k|\right).
 \label{eq:implementation}
\end{equation}
The triangle inequality bounds the received perturbation $N+\sum_k(\widetilde\enc-\enc)(s_k)+\sum_k d_k\widetilde\enc(s_k)$. If $|d_k|\le\delta$ and $\bar\lambda_\enc=\Lambda_\sq(\enc)/(K^2P)$, then
$$
 |\widehat f(Y)-f(\mathbf s)|
 \le\frac{|N|/(K\sqrt P)+\epsilon+\delta}
 {\sqrt{\bar\lambda_\enc}}.
$$
Increasing $K$ suppresses the normalized additive-noise term but not the bounds on coherently accumulated implementation and gain errors.

For a complex quantizer with error at most $r\sqrt P$, $0<r<1$, transmit
$$
 \widetilde\enc(s)=Q\bigl((1-r)\enc(s)\bigr).
$$
Then $|\widetilde\enc(s)|\le\sqrt P$ and $\epsilon\le2r$. Uniform nearest-level rounding with $B\ge2$ bits per component on $[-\sqrt P,\sqrt P]$ gives $r=\sqrt2/(2^B-1)<1$. Equation~\eqref{eq:implementation} bounds the resulting deterministic computation error. Residual amplitude and phase errors $|\alpha_k|\le\alpha_0$ and $|\theta_k|\le\theta_0\le\pi$, with $1+d_k=(1+\alpha_k)e^{i\theta_k}$, give $\delta\le\alpha_0+2\sin(\theta_0/2)$.

Retaining a fraction $\eta$ of optimal squared separation increases the disturbance factor in \eqref{eq:implementation} by at most $\eta^{-1/2}$: a $90\%$ retention limits the increase to $5.41\%$.

\section{High-SNR Design}\label{sec:pareto}
Worst-case separation protects every input. We now use the input distribution to reduce the asymptotic MSE coefficient while retaining a prescribed uniform guarantee. As defined in Section~\ref{sec:model}, we take $N_0\downarrow0$ with $P$, $K$, the encoder, and the input distribution fixed, so the per-transmitter peak SNR $P/N_0$ tends to infinity.

Throughout this section we use the hypotheses of Theorem~\ref{thm:sdp} and the strip decoder \eqref{eq:stripdecoder}. The input vector has a fixed joint distribution $\pi$, independent of the noise; its coordinates need not be independent. We assume $|m|<1$ almost surely and define
\begin{equation}
 w(m)=\begin{pmatrix}-\psi'(m)\\1\end{pmatrix},\qquad
 Q_\pi=\mathbb E_\pi[w(m)w(m)^\T]\succ0.
 \label{eq:prior-matrix}
\end{equation}
The matrix $Q_\pi\in\mathbb S^2$ averages the sensitivity of $F(m,t)=t-\psi(m)$ to perturbations of the recovered statistics $(m,t)$. Positive definiteness is equivalent to $\operatorname{Var}_\pi(\psi'(m))>0$. Bounded endpoint slopes make $Q_\pi$ finite. We allow scalar output clipping only when the true value $f(\mathbf s)$ lies strictly inside the chosen clipping interval almost surely.

\begin{proposition}\label{prop:general-small-noise}
For every continuous encoder $\enc:I\to\R^2$ satisfying the identifiability condition \eqref{eq:exactness} for $f_{\phi,\psi}:I^K\to\R$ and the peak constraint \eqref{eq:peak-model}, let $G=\mathsf U^\T\mathsf U\in\mathbb S_{++}^2$ be the Gram matrix of its coefficient vectors $u,v$. If $N\sim\mathcal{CN}(0,N_0)$, then the strip decoder $\widehat f:\R^2\to\R$ satisfies
\begin{equation}
 \mathcal J_\pi(G):=\lim_{N_0\downarrow0}
 \frac{\mathbb E_{\pi,N}|\widehat f(Y)-f(\mathbf s)|^2}{N_0}
 =\frac{\operatorname{tr}(Q_\pi G^{-1})}{2K^2}.
 \label{eq:general-small-noise}
\end{equation}
\end{proposition}
Thus $\operatorname{MSE}=N_0\mathcal J_\pi(G)+o(N_0)$: the coefficient $\mathcal J_\pi(G)$ determines the first-order MSE at high SNR.
\begin{proof}
Theorem~\ref{thm:rigidity} forces an invertible coefficient matrix $\mathsf U=[u\ v]\in\R^{2\times2}$. At every noiseless point with $|m|<1$, strip projection is locally inactive and the received-coordinate gradient is $\mathsf U^{-\T}w(m)/K$. Any permitted output clipping is also locally inactive. Couple the noises as $N=\sqrt{N_0}Z$, with $\mathbb E[ZZ^\T]=I_2/2$ in real coordinates. Differentiability gives the pointwise limit of the squared error divided by $N_0$. Theorem~\ref{thm:decoder} bounds this quotient by $L_\enc^2|Z|^2$, an integrable random variable. Dominated convergence yields \eqref{eq:general-small-noise}.
\end{proof}

We can now choose how much uniform robustness to retain while minimizing the asymptotic MSE coefficient. Let $\lambda_\star$ denote the optimum in \eqref{eq:sdp}. For a required separation floor $0\le\gamma\le\lambda_\star$, define $\mathcal V_\pi(\gamma)$ as the smallest coefficient \eqref{eq:general-small-noise} among all continuous encoders permitting identifiable computation of $f_{\phi,\psi}$ under the peak constraint and $\Lambda_\sq\ge K^2\gamma$. This design guarantees uniform amplification at most $1/(K\sqrt\gamma)$ when $\gamma>0$.

\begin{theorem}[Robustness--accuracy frontier]\label{thm:pareto}
Using the enclosing-ellipse class $\mathcal C_P$ in \eqref{eq:ellipse-class}, for $\gamma>0$ we obtain $\mathcal V_\pi(\gamma)$ from the five-variable convex program
\begin{equation}
\begin{aligned}
 \text{minimize}\quad &\frac{\operatorname{tr}(Q_\pi X)}{2K^2}\\
 \text{over}\quad &X\in\mathbb S^2,\quad a\in\R^2,\\
 \text{subject to}\quad
 &(X,a)\in\mathcal C_P,\\
 &w_\sigma^\T Xw_\sigma\le\gamma^{-1}
       \quad(\sigma\in\{-,+\}).
\end{aligned}
 \label{eq:pareto-sdp}
\end{equation}
At $\gamma=0$, we omit the last two inequalities. An optimum exists, has a unique inverse Gram matrix $X=G^{-1}\in\mathbb S_{++}^2$, and yields a planar encoder permitting identifiable computation of $f_{\phi,\psi}$ via
$$
 \mathsf U^\T\mathsf U=X^{-1},\qquad
 \enc(s)=\mathsf U(g(s)-a).
$$
The value is finite, continuous, convex, and nondecreasing on $[0,\lambda_\star]$.

Let $X_0$ minimize \eqref{eq:pareto-sdp} at zero floor, and set
$$
 \gamma_c=\left(\max_{\sigma\in\{-,+\}}
 w_\sigma^\T X_0w_\sigma\right)^{-1}.
$$
The frontier is constant on $[0,\gamma_c]$. The encoder minimizing the asymptotic MSE coefficient already meets every separation requirement in this interval. If $\gamma_c<\lambda_\star$, it is strictly increasing thereafter, and every optimal encoder on that branch satisfies $\Lambda_\sq=K^2\gamma$.
\end{theorem}
The proof appears in Appendix~\ref{app:general-frontier-proof}.

This theorem optimizes encoders for a fixed decoder architecture and the asymptotic MSE coefficient. It does not optimize finite-SNR risk over arbitrary estimators, nor interchange the limit $N_0\downarrow0$ with a noise-dependent encoder optimization. The objective and endpoint constraints in \eqref{eq:pareto-sdp} are linear, but its peak condition remains a continuum of matrix inequalities; a finite-grid solver requires a separate certificate.

\subsection{Quadratic Frontier}
For the quadratic preprocessing function $\phi(s)=s^2$, symmetry reduces the matrix tradeoff to one encoder coefficient. The distinction between the two objectives then appears directly as a worst-case slope and a root-mean-square slope. For $\phi(s)=s^2$, even $\psi$, and a symmetric distribution of $m$, set
$$
 M=|\psi'(1)|,\qquad M_{\rm eff}=\sqrt{\mathbb E_\pi[\psi'(m)^2]}.
$$
Our assumptions imply $0<M_{\rm eff}<M$, and $Q_\pi=\operatorname{diag}(M_{\rm eff}^2,1)$. For $x>0$, define the decreasing function
$$
 b(x)=\frac{2}{1-(x/2)^2+(x/2)\sqrt{(x/2)^2+2}},
$$
and set $b_M=b(M)$, $b_{\rm eff}=b(M_{\rm eff})$. For $1\le b<2$, write
$$
 R_M(b)=\left(\frac{M^2}{b(2-b)}+\frac1{b^2}\right)^{-1}.
$$
Its unique maximum occurs at $b_M$, and $\lambda_\star=P R_M(b_M)$. For $0<r\le R_M(b_M)$, let $b_+(r)$ be the unique solution of $R_M(b)=r$ on $[b_M,2)$, and define $b_+(0)=2$. Equivalently, the selected root satisfies
$$b^3-2b^2+r(M^2-1)b+2r=0.$$

\begin{corollary}\label{cor:quadratic-pareto}
For every $0\le\gamma\le\lambda_\star$, define
$b_\gamma=\min\{b_{\rm eff},b_+(\gamma/P)\}.$
An encoder attaining the global frontier over all continuous encoders permitting identifiable computation of $f_{\phi,\psi}$ under the peak constraint is
\begin{equation}
 \enc_\gamma(s)=\sqrt P\left[\sqrt{b_\gamma(2-b_\gamma)}\,s
       +i(b_\gamma s^2-1)\right],
 \label{eq:quadratic-pareto-map}
\end{equation}
and its optimal asymptotic MSE coefficient is
\begin{equation}
 \mathcal V_\pi(\gamma)=\frac1{2K^2P}
 \left[\frac{M_{\rm eff}^2}{b_\gamma(2-b_\gamma)}
       +\frac1{b_\gamma^2}\right].
 \label{eq:quadratic-pareto-value}
\end{equation}
The constant portion ends at $\gamma_c=P R_M(b_{\rm eff})$.
\end{corollary}
The proof appears in Appendix~\ref{app:quadratic-frontier-proof}. Increasing $b_\gamma$ allocates more of the available scale to the second moment $t=q$ and less to the sample mean $m$. The separation floor limits this reallocation because inputs near the endpoint means remain sensitive to mean-estimation errors.

\begin{figure*}[t]
 \centering\paperfigure{figures/tikz/pareto_frontier}
 \caption{Variance design with a prescribed separation floor. Left: reduction of the optimal asymptotic MSE coefficient relative to the separation-optimal encoder; markers identify where the required separation floor becomes active. Right: the preprocessing curve $g(s)=(s,s^2)^\T$ and containing ellipses for full robustness, $90\%$ retention, and the uniform-$K=16$ prior optimum. The ellipses express the peak constraint in the coordinates of $g(s)$; they are not noise confidence regions.}
 \label{fig:pareto}
\end{figure*}

Figure~\ref{fig:pareto} relates this frontier to the enclosing ellipses of Section~\ref{sec:inverse-synthesis}. For variance and independent uniform sources, $M=2$ and $M_{\rm eff}=2/\sqrt{3K}$. The choice $\gamma=\eta\lambda_\star$, $0\le\eta\le1$, therefore specifies the entire tradeoff analytically through a single scalar root. It interpolates between the encoder minimizing the asymptotic MSE coefficient and the separation-optimal encoder while making the retained uniform guarantee explicit.

\section{Certified Design}\label{sec:certified}
For general $\phi\in C^2(I)$, we certify the continuum peak constraints in \eqref{eq:sdp} using finite semidefinite optimization, an interpolation bound, and verified primal--dual witnesses. This additional smoothness serves certification, not encoder classification. Numerical implementation requires verified enclosures of $\phi$, $\psi'(\pm1)$, and the derivative bounds below.

\begin{theorem}\label{thm:certifiedgrid}
Assume the hypotheses of Theorem~\ref{thm:sdp} and $\phi\in C^2(I;\R)$. Let $\mathcal T_h\subset I$ be a finite grid containing both endpoints, with maximum adjacent spacing $h>0$, and containing three fixed inputs whose vectors $z(s)=(1,s,\phi(s))^\T$ span $\R^3$. Let $U_h$ denote the optimum obtained by imposing the peak constraints only on $\mathcal T_h$, and let $(H_h,U_h)\in\mathbb S_+^3\times[0,\infty)$ be an optimizer. Set
$$
 Z_j=\max_{s\in I}\|z^{(j)}(s)\|_2,\qquad j=0,1,2,
$$
where $z^{(0)}=z$, and define
$$
 M_h=2\|H_h\|_2(Z_1^2+Z_0Z_2),\qquad
 \rho_h=\frac{P}{P+M_hh^2/8}.
$$
Then $(\rho_h H_h,\rho_h U_h)$ satisfies every constraint of \eqref{eq:sdp}, and
\begin{equation}
 \rho_h U_h\le\lambda_\star\le U_h,
 \qquad U_h-\rho_h U_h=O(h^2)\quad(h\downarrow0).
 \label{eq:certifiedgrid}
\end{equation}
\end{theorem}
The interpolation remainder bounds the peak between adjacent grid points; scaling then restores the continuous peak constraint. Appendix~\ref{app:grid-proof} gives the proof.

Define $q_H(s)=z(s)^\T Hz(s)$. Suppose a candidate $(H,\lambda)$ with $\lambda>0$ satisfies the semidefinite constraints, a verified number $b_h$ bounds $q_H$ at all grid points, and $M$ bounds $|q_H''|$. Then
$$
 L=\frac{P\lambda}{b_h+Mh^2/8}
$$
is an achieved lower bound, provided the denominator is positive: scale $H$ by the same factor and apply the planar reconstruction in Theorem~\ref{thm:sdp}. A tighter derivative bound than the norm bound above is equally valid.

For a global upper bound, let $\alpha_i\ge0$ be weights at any finitely many inputs $t_i\in I$. Define $\overline D=\operatorname{diag}(0,D_-+D_+)$ for $D_-,D_+\in\mathbb S_+^2$ and require
\begin{align}
 Q&=\sum_i\alpha_i z(t_i)z(t_i)^\T-\overline D\suc0,\label{eq:dualwitness}\\
 w_-^\T D_-w_-+w_+^\T D_+w_+&\ge1.\nonumber
\end{align}
Here $\overline D$ has a scalar top-left block. Every feasible pair in \eqref{eq:sdp} obeys
$$
 P\sum_i\alpha_i\ge\sum_i\alpha_i z(t_i)^\T Hz(t_i)
 \ge\sum_{\sigma\in\{-,+\}}\operatorname{tr}(D_\sigma G)
 \ge\lambda.
$$
Thus $U=P\sum_i\alpha_i$ is a global upper bound. The first omitted term is $\operatorname{tr}(QH)\ge0$, and the last inequality uses both score constraints and the normalization in \eqref{eq:dualwitness}. These witnesses yield the stopping rule $(U-L)/U\le\varepsilon$ for $U>0$. Exact grid solutions converge by Theorem~\ref{thm:certifiedgrid}; in computation, primal and dual verification errors must also decrease to obtain arbitrarily small certified gaps.

\subsection{Exponential Preprocessing}
For $\phi(s)=e^s$ and $\psi(m)=e^m$, we certify a design for
$$
 f:I^K\to\R,\qquad f(\mathbf s)=\frac1K\sum_{k=1}^K e^{s_k}-e^{K^{-1}\sum_{k=1}^K s_k}.
$$
We normalize $P=1$; multiplication of the encoder by $\sqrt P$ restores arbitrary peak power. SciPy's sequential quadratic programming routine supplies a candidate for the finite convex problem on 401 grid points. We then round and repair its primal and dual matrices and verify the resulting witnesses independently of solver tolerances.

All certificate arithmetic uses rational intervals. At each rational input $|s|\le1$, we enclose $e^s$ by its degree-40 Taylor polynomial with remainder bounded by $3/41!$. We certify positive definiteness through strictly positive lower bounds on the leading principal minors. For the primal witness, a 1001-point grid gives $b_h\le1.000000093894$, and direct differentiation gives $|q_H''|\le117$. Hence $h=1/500$ and the verified continuous peak bound before scaling is $1.000058593894$. The dual witness uses the three candidate contact inputs $-1,159/200,1$, together with positive repair weights at $-1,0,1$.

The exact rational witnesses establish
\begin{equation}
 0.419645089928\le
 \frac{\Lambda_{\sq}^\star}{K^2P}
 \le0.419670224438.
 \label{eq:expcertificate}
\end{equation}
Thus the feasible encoder attains at least $99.994\%$ of the optimal squared-separation score, with a certified relative gap below $0.006\%$.

For orientation, the reconstructed curve has approximate coordinates
\begin{align*}
 \frac{\operatorname{Re}\enc(s)}{\sqrt P}
 &\approx1.8321854+2.1939971s-1.1944545e^s,\\
 \frac{\operatorname{Im}\enc(s)}{\sqrt P}
 &\approx-0.7900289+0.5211418e^s.
\end{align*}
The certificate uses the stored rational Gram matrix and exact reconstruction, not these rounded display coefficients. Figure~\ref{fig:certified} shows the resulting curve and the certified gap as we refine only the verification grid while holding the candidate and dual witness fixed. The limiting gap of that fixed-candidate experiment also includes the candidate and witness repair errors. 

\begin{figure*}[t]
 \centering\paperfigure{figures/tikz/certified_synthesis}
 \caption{Certified design for exponential preprocessing, $\phi(s)=e^s$. Left: the reconstructed feasible encoder and the unit peak boundary; markers identify the candidate contact inputs used in the certificate. Right: the global relative-gap bound as the verification grid increases from 51 to 1001 points, with a fixed candidate and dual upper bound. The repaired encoder need not meet the peak boundary exactly.}
 \label{fig:certified}
\end{figure*}

\section{Numerical Results}\label{sec:evidence}
We now examine how the design guarantees translate into numerical performance. We first check the implemented curves, then evaluate the robustness--MSE tradeoff, and finally vary the noise power to assess the range in which the small-noise predictions describe the measured error.

\subsection{Analytic Checks}
We compare the scores of the curves in Fig.~\ref{fig:curves}. The analytic proofs establish global optimality; the numerical checks validate coefficients and implementations.

\begin{table}[t]\caption{Normalized squared scores and sharp amplification. Here $\bar\lambda=\Lambda_\sq/(K^2P)$ and $\ell=K\sqrt P L_\enc=1/\sqrt{\bar\lambda}$. Baselines are defined in Fig.~\ref{fig:curves}.}\label{tab:values}
\centering\small\begin{tabular}{@{}lrrr@{}}\toprule
Desired function & Baseline $\bar\lambda$ & Optimal $\bar\lambda$ & Optimal $\ell$\\\midrule
$q-m^2$ & $0.200000$ & $0.206267$ & $2.201835$\\
$t_4-m^4$ & $0.047059$ & $0.058534$ & $4.133308$\\
$mq$ & $0.500000$ & $0.592593$ & $1.299038$\\\bottomrule
\end{tabular}\end{table}


The quadratic baseline is feasible by the peak identity with $b=1$. For the quartic baseline, $x=s^2\in[0,1]$ gives
$$|\enc_{\rm base}(s)|^2/P=\tfrac45[x+(x^2-1/2)^2]\le1,$$
since $x\le1$ and $(x^2-1/2)^2\le1/4$. Its exact score is $(4/85)K^2P$. Table~\ref{tab:values} compares these specified reference designs with the certified optima, without implying MSE gains.

The numerical reproduction package includes the checkers, rational certificates, simulation seeds, and recorded results used in this section.\footnote{\href{https://github.com/SaeedRazavikia/Optimal-Continuous-Analog-Encoders-TIT-/blob/7a513ed38147a73f109f9b184dafba59c17ee89b/REPRODUCIBILITY.md}{Reproduction instructions and numerical source, commit \texttt{7a513ed}.} The accompanying source archive also includes these files.} The structural checker verifies the implemented coefficients, quadratic stationarity, quartic peak and multiplier identities, asymmetric mean penalties, shrinking source perturbations, and Gram rank reduction. A separate decoder checker compares the strip formula for nonorthogonal directions with received-plane quadratic programming and verifies exactness and disturbance bounds across several values of $K$ and $P$. These checks pass, but their purpose is implementation verification. The full-interval analytic identities and the rational certificate in Section~\ref{sec:certified} establish feasibility and global bounds independently of sampling.



\FloatBarrier
\subsection{Robustness--MSE Tradeoff}
\label{sec:pareto-evidence}
We evaluate Corollary~\ref{cor:quadratic-pareto} for variance, $P=1$, $K\in\{4,16\}$, and independent uniform or endpoint-Beta sources, defined by $s=2X-1$ with $X\sim\operatorname{Beta}(0.3,0.3)$. We write $\eta=\gamma/\lambda_\star$ for the required fraction of optimal squared separation. Figure~\ref{fig:pareto} shows the exact small-noise frontier and its enclosing-ellipse geometry.



For uniform sources, the prior optimum retains $83.91\%$ of separation at $K=4$ and $60.91\%$ at $K=16$. A $90\%$ floor is active in both cases, giving $a=0.9038538$, $b=1.4278414$, and $c=-1$. The small-noise reductions are $17.68\%$ and $29.07\%$, respectively, while the sharp disturbance factor increases by $5.41\%$. For endpoint-Beta sources at $K=4$, the prior optimum retains $91.67\%$, so the $90\%$ floor is inactive.

We evaluate eight prescribed retentions per scenario at $P/N_0=20$ dB using $200{,}000$ independent source/noise draws, shared across encoders. Table~\ref{tab:pareto} compares predicted and measured gains at the $90\%$ floor, with nominal pointwise $95\%$ confidence intervals from the paired delta method. These finite-noise measurements assess the specified designs; the theorem concerns the limiting coefficient.

\begin{table}[t]
\caption{MSE reductions (\%) from the separation-optimal variance encoder at a $90\%$ separation floor. U: uniform; B: endpoint Beta.}
\label{tab:pareto}
\centering\small
\begin{tabular}{@{}ccrr@{}}\toprule
Prior & $K$ & Small-noise limit & $20$ dB, paired $95\%$ interval\\\midrule
U & 4 & 17.68 & $17.99\ [17.84,18.14]$\\
U & 16 & 29.07 & $29.10\ [29.02,29.19]$\\
B & 4 & 9.79 & $10.85\ [10.70,11.00]$\\
B & 16 & 25.09 & $25.12\ [25.00,25.24]$\\\bottomrule
\end{tabular}
\end{table}

Twelve comparisons with a separately parameterized convex problem agree within $1.04\times10^{-9}$ relative objective error. Two runs report line-search termination despite residuals within $2.41\times10^{-10}$; these are consistency checks, not certificates. A second implementation reproduces the two uniform $90\%$ points, exponential rational certificate, sixteen coefficient/moment cases, and nondiagonal decoder comparison. Both implementations use the same environment and do not constitute external reproduction.

We next vary the noise power to assess clipping and nonlinear receiver effects beyond this regime.

\FloatBarrier
\subsection{Finite-Noise Performance}\label{sec:finite-noise}
We next evaluate the same encoders with the projection decoder of Section~\ref{sec:decoding}, followed by clipping to $[0,1]$. This comparison separates uniform disturbance robustness from average performance under a specified input distribution. For $p\in\{2,4\}$, write
$$
 f_p=t_p-m^p,\qquad
 \enc(s)=\sqrt P\,[as+i(bs^p+c)],\quad a,b>0.
$$
Let $N=\sqrt{N_0}Z$, where $Z$ is circular Gaussian with $\mathbb E|Z|^2=1$, independent of the sources. Thus each real noise component has variance $N_0/2$.

For independent sources with a common continuous density on $(-1,1)$, the target lies strictly inside $(0,1)$ almost surely. Proposition~\ref{prop:general-small-noise} therefore specializes to
\begin{equation}
 \lim_{N_0\downarrow0}\frac{\operatorname{MSE}(\enc,N_0)}{N_0}
 =\frac{1}{2K^2P}\left(\frac{1}{b^2}
 +\frac{p^2\mathbb E[m^{2p-2}]}{a^2}\right).
 \label{eq:small-noise-mse}
\end{equation}


The constrained theorem in Section~\ref{sec:pareto} interpolates between these objectives. The worst-case criterion uses the endpoint slope $M=p$, whereas \eqref{eq:small-noise-mse} uses $M_{\rm eff}=p\sqrt{\mathbb E[m^{2p-2}]}$. We construct a prior-aware comparison curve by inserting $M_{\rm eff}$ into the quadratic or quartic coefficient formulas. This choice minimizes the leading small-noise coefficient within the stated orthogonal curve family. We evaluate its moments analytically, without fitting simulation data, and use the same decoder for all designs. This comparison does not optimize finite-noise Bayesian risk.




We set $P=1$, $K\in\{4,16\}$, and draw independent uniform or endpoint-Beta sources. For each input distribution and transmitter count, we generate $200{,}000$ input vectors and independent standardized Gaussian disturbances. We reuse these draws across designs, targets, and the nine SNR values from $-10$ to $30$ dB in $5$-dB steps. Here SNR means the per-transmitter peak ratio $P/N_0$, not an aggregate received-power ratio. Paired standard errors account for the shared draws; intervals are pointwise, without a multiple-comparison adjustment.

\begin{table}[t]
\caption{MSE at $P/N_0=20$ dB, in units of $10^{-4}$. U: uniform; B: endpoint Beta. Robust: separation-optimal curve; tuned: prior-aware small-noise curve.}
\label{tab:finite-noise}
\centering\small
\begin{tabular}{@{}cccrrr@{}}\toprule
$p$ & Prior & $K$ & Baseline & Robust & Tuned\\\midrule
2 & U & 4  & 4.055 & 3.318 & 2.684\\
2 & B & 4  & 4.787 & 4.075 & 3.634\\
2 & U & 16 & 0.2115 & 0.1630 & 0.09922\\
2 & B & 16 & 0.2252 & 0.1772 & 0.1232\\
4 & U & 4  & 4.137 & 3.341 & 1.652\\
4 & B & 4  & 5.627 & 4.548 & 3.338\\
4 & U & 16 & 0.2450 & 0.1970 & 0.05521\\
4 & B & 16 & 0.2471 & 0.1987 & 0.06590\\\bottomrule
\end{tabular}
\end{table}

Across the 72 tested target/prior/$K$/SNR combinations, the separation-optimal curves reduce MSE by $6.12\%$--$22.92\%$ relative to the specified baselines. However, the prior-aware curves further reduce MSE by $5.69\%$--$72.01\%$ relative to the separation-optimal curves. Figure~\ref{fig:finite-noise} and Table~\ref{tab:finite-noise} display both outcomes. These average-performance gains sacrifice worst-case robustness: for the quartic gap with uniform sources and $K=16$, prior-aware tuning reduces $\Lambda_\sq/(K^2P)$ from $0.05853$ to $0.007273$. At low SNR, even the known prior-mean estimator can perform better: for uniform sources, $K=4$, variance, and $-10$ dB, its MSE is $0.02294$, compared with $0.1032$ for the separation-optimal curve and implemented decoder. These results quantify the operating behavior of specified designs; they establish neither universal MSE dominance nor Bayesian optimality. At $30$ dB, the endpoint-Beta, $K=4$ normalized errors still lie $3.49\%$--$7.19\%$ below their small-noise limits, beyond Monte Carlo uncertainty; boundary clipping remains relevant at this finite noise level. The code records every tested condition, paired uncertainty, analytic prior moments, seeds, and decoder cross-checks.


\FloatBarrier
\section{Conclusions}\label{sec:boundaries}
We classify exact shared continuous planar encoders for the desired functions $f_{\phi,\psi}(\mathbf s)=t-\psi(m)$ and optimize their worst-case conditioning under pointwise peak power. We turn this structural result into a sharp uniform disturbance guarantee by combining the exact score, convex encoder design, and projection decoder. The inverse-metric description also yields the global robustness-constrained small-noise frontier for the specified decoder. Its quadratic specialization quantifies the cost of adapting to an input distribution, while certified discretization supplies a stopping rule for smooth local preprocessing functions $\phi$ beyond the explicit cases.

The assumptions delimit the contribution. We use one synchronized complex channel use, unit effective channel gains, the same encoder at every transmitter, and at least four transmitters. Unequal uncompensated gains, user-specific encoders, average-power constraints, and unrestricted $K=3$ classification require different arguments. The desired function $f:I^K\to\R$ is fixed before design; the result does not cover arbitrary symmetric functions. The extension in Section~\ref{sec:extension} requires common-completion identifiability for encoder classification, Lipschitz continuity for a positive design optimum, and separate endpoint conditions for the explicit quadratic formula. Universal sum representations have established dimension limitations~\cite{wagstaff19,wagstaff22}.

For a fixed encoder, the score determines the smallest uniform disturbance-amplification constant over all exact decoders. The constrained frontier minimizes the leading MSE coefficient of the specified strip decoder; finite-noise Bayesian optimization over arbitrary estimators remains a separate question. Section~\ref{sec:implementation} bounds residual mismatch without solving optimal design for unequal channels. Appendix~\ref{app:boundaries} gives precise boundaries for criterion transfer, poor conditioning, and finite alphabets. Appendix~\ref{app:approx-rigidity} proves qualitative stability under common regularity bounds, while an explicit error-to-geometry rate remains open. These restrictions identify concrete directions for extending the present classification and design theorem.

\appendices
\numberwithin{theorem}{section}
\section{Proof of Theorem~\ref{thm:rigidity}}\label{app:rigidity-proof}
We separate the proof into three steps because continuity alone does not permit us to differentiate the encoder. A fourth transmitter first turns equality of encoded triple sums into equality of both moments. We then use topology to obtain local dependence on those moments and regularization to justify differentiation. The resulting derivative identities determine the encoder coefficients.

\begin{lemma}
\label{lem:collision}
Let $K\ge4$, $I=[-1,1]$, $\phi:I\to\R$, and $\psi:I\to\R$, where $\psi$ is strictly convex or strictly concave. Suppose $f_{\phi,\psi}:I^K\to\R$ in \eqref{eq:target} is identifiable through $\enc:I\to\R^2$. Then, for every $\mathbf a,\mathbf b\in(-1,1)^3$,
\begin{align*}
 &\sum_{j=1}^3\enc(a_j)=\sum_{j=1}^3\enc(b_j)\\
 &\quad\Longrightarrow\quad
 \begin{pmatrix}\sum_{j=1}^3 a_j\\\sum_{j=1}^3\phi(a_j)\end{pmatrix}
 =\begin{pmatrix}\sum_{j=1}^3 b_j\\\sum_{j=1}^3\phi(b_j)\end{pmatrix}.
\end{align*}
\end{lemma}
\begin{proof}
Let the triple statistics be $(S,T)$ and $(\widetilde S,\widetilde T)$. Complete both triples with the same fourth input $w$ and $K-4$ zeros. The identifiability condition gives
$$
 T-\widetilde T
 =K\left[\psi\!\left(\frac{S+w}{K}\right)
 -\psi\!\left(\frac{\widetilde S+w}{K}\right)\right].
$$
Use $w=0$ and $w=1/2$. With $a=S/K$, $\widetilde a=\widetilde S/K$, and $\delta=1/(2K)>0$, subtraction yields
$$
 \psi(a+\delta)-\psi(a)
 =\psi(\widetilde a+\delta)-\psi(\widetilde a).
$$
All four arguments lie in $(-1,1)$. If $a<\widetilde a$, the points $a+\delta$ and $\widetilde a$ lie strictly between $a$ and $\widetilde a+\delta$ with complementary convex weights. Strict convexity therefore gives
$$
 \psi(a+\delta)+\psi(\widetilde a)
 <\psi(a)+\psi(\widetilde a+\delta),
$$
contradicting the equal increments. Strict concavity reverses this inequality, and exchanging the triples excludes $a>\widetilde a$. Thus $S=\widetilde S$, and the same identifiability condition at $w=0$ gives $T=\widetilde T$. This proof also permits endpoint inputs and does not require differentiability of $\psi$.
\end{proof}

The collision lemma ensures that the received sum determines the moments. We must also show that it cannot vary locally while those moments remain fixed. The next lemma supplies this converse dependence using invariance of domain: a continuous injection from an open subset of $\R^2$ into $\R^2$ has open image~\cite[Theorem 2B.3]{hatcher}.

\begin{lemma}
\label{lem:topology}
Let $B\subset\R^2$ be open, let $J\subset\R$ be an interval, and let $R:B\times J\to\R^2$ be continuous. Suppose there is a map $q:R(B\times J)\to B$ such that $q(R(p,t))=p$ for every $(p,t)\in B\times J$. Then $R(p,t)$ is independent of $t$ for each fixed $p$.
\end{lemma}
\begin{proof}
For fixed $t_0$, the map $p\mapsto R(p,t_0)$ is a continuous injection, so $R(B,t_0)$ is open. Fix $p$. For $t$ sufficiently close to $t_0$, continuity gives $R(p,t)\in R(B,t_0)$. Thus $R(p,t)=R(p',t_0)$ for some $p'\in B$. Applying $q$ yields $p'=p$, so $R(p,t)=R(p,t_0)$. The function $t\mapsto R(p,t)$ is locally constant, and connectedness of $J$ makes it constant. This argument does not require differentiability or continuity of $q$.
\end{proof}

The final lemma adapts restricted Cauchy--Pexider rigidity~\cite[Theorem 1.6 and Sec. 5]{charalambides2013} to a one-way collision implication and a nonaffine local preprocessing function $\phi$ of class $C^1$ that may contain affine subintervals.

\begin{lemma}
\label{lem:affine}
Let $\phi:I\to\R$ be nonaffine and belong to $C(I)\cap C^1((-1,1))$, and let $\enc:I\to\R^2$ be continuous. Suppose equality of encoded sums for any two interior triples implies equality of their $S$ and $T$ statistics. Then $\enc$ has form \eqref{eq:affine} with $c,u,v\in\R^2$ and linearly independent $u,v$.
\end{lemma}
\begin{proof}
\emph{Step 1: Local factorization through the moments.}
Choose interior $y_0,z_0$ with $\phi'(y_0)\ne\phi'(z_0)$; such a pair exists because $\phi$ is nonaffine. At any $(x_0,y_0,z_0)$, the map
$$
 (x,y,z)\longmapsto
 \bigl(x,\;x+y+z,\;\phi(x)+\phi(y)+\phi(z)\bigr)
$$
has Jacobian determinant $\phi'(z_0)-\phi'(y_0)\ne0$ and is a local $C^1$ diffeomorphism. Restrict its image to a product $J\times B$, where $x_0\in J$ and $B\subset\R^2$ is open. In these coordinates, the triple received sum is a continuous $R(p,x)$. The collision hypothesis makes the assignment $q(R(p,x))=p$ well defined even when distinct points of the chart have the same received sum. Lemma~\ref{lem:topology} shows that $R$ is independent of $x$. In particular, $H(p)=R(p,x_0)$ is continuous. On a sufficiently small source box whose image lies in $J\times B$, we have
\begin{equation}
 \enc(x)+\enc(y)+\enc(z)
 =H\bigl(x+y+z,\phi(x)+\phi(y)+\phi(z)\bigr)
 \label{eq:localfactor}
\end{equation}
for a continuous $H$ defined on $B$.

\emph{Step 2: Automatic interior differentiability.}
The pair map $(y,z)\mapsto w=(y+z,\phi(y)+\phi(z))$ is a local $C^1$ diffeomorphism near $(y_0,z_0)$. Restrict this map to the pair-coordinate part of the source box from Step~1. Choose a smooth compactly supported density $\rho(w)$ of integral one inside its open image, and denote its inverse by $(y(w),z(w))$. Integrating \eqref{eq:localfactor} with respect to $\rho(w)\,dw$ gives, for $x$ in the common source interval,
$$
 \enc(x)=\int H\bigl((x,\phi(x))+w\bigr)\rho(w)\,dw-C_0,
$$
where $C_0=\int[\enc(y(w))+\enc(z(w))]\rho(w)\,dw$ is independent of $x$. Choose the source interval $X$ and support of $\rho$ small enough that $\{(x,\phi(x))+w:x\in\overline X,\ w\in\operatorname{supp}\rho\}$ is compactly contained in $B$. Choose a smooth compactly supported cutoff in $B$ equal to one near this set, multiply $H$ by it, and extend the product by zero. The integral is then a smooth convolution as a function of the two-dimensional argument. Its composition with the $C^1$ curve $x\mapsto(x,\phi(x))$ is $C^1$. Since $x_0$ was arbitrary, $\enc$ is $C^1$ throughout the interval interior. No regularity of the original decoder was assumed.

\emph{Step 3: Affine dependence on $s$ and $\phi(s)$.}
Consider a real component $e$ of the now differentiable encoder. Equation~\eqref{eq:localfactor} makes $e(x)+e(y)+e(z)$ constant on each local moment fiber. The two-moment map has rank two whenever $\phi'(y)\ne\phi'(z)$, so the implicit-function theorem identifies its differential kernel with the tangent space of that fiber. The differential of the encoded sum annihilates this kernel. Hence
$$
 (e'(x),e'(y),e'(z))
 \in\operatorname{span}\{(1,1,1),(\phi'(x),\phi'(y),\phi'(z))\}.
$$
Fix $y=y_0,z=z_0$. Their entries give
$$
\beta=\frac{e'(y_0)-e'(z_0)}{\phi'(y_0)-\phi'(z_0)},
\qquad \alpha=e'(y_0)-\beta\phi'(y_0).
$$
These constants do not depend on $x$. Since Step~1 applies at every $(x,y_0,z_0)$, we obtain $e'(x)=\alpha+\beta\phi'(x)$ throughout the connected interval $(-1,1)$. Integration, followed by endpoint continuity, proves \eqref{eq:affine} componentwise.

Finally, the triple moment map has an open image locally near $(x_0,y_0,z_0)$. If $u,v$ were dependent, the linear map $(S,T)\mapsto uS+vT$ could not be injective on that open set. This would contradict the collision hypothesis. Thus $u,v$ are independent.
\end{proof}

\begin{proof}[Proof of Theorem~\ref{thm:rigidity}]
Necessity follows from Lemmas~\ref{lem:collision} and~\ref{lem:affine}. For sufficiency, the receiver solves the invertible real system $r_\enc-Kc=uS+vT$ and evaluates $T/K-\psi(S/K)$.
\end{proof}

\section{Proof of Theorem~\ref{thm:sdp}}\label{app:design-proof}
We first compute the exact score under the hypotheses of Section~\ref{sec:synthesis}. The key issue is sharpness: a gradient bound must be approached by actual input configurations. We then show that the matrix optimization preserves both this score and planar realizability.

\begin{proposition}
\label{prop:score}
Under the hypotheses of Section~\ref{sec:synthesis}, let $\enc:I\to\R^2$ have the form \eqref{eq:affine} with $c,u,v\in\R^2$ and linearly independent $u,v$. Its operational score for $f_{\phi,\psi}:I^K\to\R$ is given by \eqref{eq:exactscore}.
The infimum is approached by pairs with received distance tending to zero.
\end{proposition}
\begin{proof}
In the moment coordinates $(S,T)$, the desired function is $F_K(S,T)=T/K-\psi(S/K)$, with gradient $K^{-1}(-\psi'(S/K),1)^\T\in\R^2$. The squared received distance is $\Delta p^\T G\Delta p$, where here $\Delta p=(\Delta S,\Delta T)^\T$. Since $\psi'$ is monotone, it ranges between its endpoint values. The quadratic function $(-q,1)G^{-1}(-q,1)^\T$ is convex in $q$, so its maximum occurs at an endpoint. Integrating the gradient along the straight segment between two moment points proves the lower bound in \eqref{eq:exactscore}. The segment need not be attainable: the formula $T/K-\psi(S/K)$ is defined on the entire strip $|S|\le K$.

For sharpness, fix an interior mean $m$ and choose two distinct nearby inputs $x,y$ with $(x+y)/2=m$; set the remaining inputs to $m$. Strict convexity or concavity of $\phi$ gives $\phi'(x)\ne\phi'(y)$. For any moment direction $d=(d_S,d_T)^\T$, choose $\alpha,\beta$ to solve
$$
 \begin{pmatrix}1&1\\\phi'(x)&\phi'(y)\end{pmatrix}
 \begin{pmatrix}\alpha\\\beta\end{pmatrix}=d.
$$
Vary only these two inputs along $x+h\alpha,y+h\beta$. The input vector remains admissible for sufficiently small $h$, and its moment derivative is $d$. With $w(m)=(-\psi'(m),1)^\T$ and $d=G^{-1}w(m)$, the derivative of the desired function along this perturbation is $w(m)^\T G^{-1}w(m)/K>0$. The separation quotient along this path therefore tends to $K^2/[w(m)^\T G^{-1}w(m)]$. Let $m$ approach the endpoint with the larger dual norm. At each mean choose a nonzero displacement small enough that the received distance is below $1/n$ and the quotient is within $1/n$ of its local limit. Continuity of the endpoint slopes gives equality and the shrinking-pair statement. This construction uses actual input perturbations and requires no uniform perturbation radius near the endpoint.
\end{proof}
\subsection{Optimality and Attainment}
With the operational score established, the remaining task is to justify the Gram formulation. In particular, a feasible $3\times3$ matrix must yield coefficients in the two-dimensional transmit plane without increasing their peak energy.
\begin{proof}[Proof of Theorem~\ref{thm:sdp}]
Every planar coefficient triple $(c,u,v)$ defines its $3\times3$ Gram matrix $H$. The peak inequalities are exactly the stated quadratic forms. For $G\succ0$, the identity
$$
 G-\lambda ww^\T\suc0
 \quad\Longleftrightarrow\quad
 \lambda\le\frac1{w^\T G^{-1}w}
$$
shows that the last two constraints enforce the score in Proposition~\ref{prop:score}. Thus \eqref{eq:sdp} upper-bounds every positive-score continuous encoder, by Theorem~\ref{thm:rigidity}.

Conversely, write $H=\left(\begin{smallmatrix}h&\ell^\T\\\ell&G\end{smallmatrix}\right)$. At every feasible point with $\lambda>0$, the two distinct endpoint slopes imply $G\succ0$: a null vector of $G$ would have to be orthogonal to both independent vectors $w_-,w_+$. Replace $h$ by $\ell^\T G^{-1}\ell$. The resulting matrix remains positive semidefinite, has rank two, preserves $G$, and decreases every peak quadratic form by the same nonnegative constant. It is therefore the Gram matrix of three vectors in $\R^2$, giving a feasible exact encoder with score at least $K^2\lambda$.

The reconstruction given in Section~\ref{sec:synthesis} realizes this reduced Gram matrix.

Finally, a sufficiently small nondegenerate moment curve is feasible and has positive score. Nonaffineness of $\phi$ supplies three inputs for which the vectors $z(s)$ span $\R^3$. If $Z$ has those three vectors $z(s)$ as columns, $Z^\T HZ\suc0$ has trace at most $3P$, giving $\|H\|_2\le3P\|Z^{-1}\|_2^2$. The last diagonal score constraint gives $0\le\lambda\le G_{22}$. The constraints are closed, and $\lambda$ is bounded by either final matrix inequality. An optimum therefore exists, is positive, and admits the rank-two reconstruction above.
\end{proof}
\subsection{Enclosing Ellipses}
\begin{corollary}\label{cor:inverse-synthesis}
Under the hypotheses of Theorem~\ref{thm:sdp}, program~\eqref{eq:inverse-synthesis} attains a positive optimum $t_\star$ and gives $\Lambda_\sq^\star=K^2/t_\star$. An optimal encoder follows from the enclosing-ellipse reconstruction in Section~\ref{sec:inverse-synthesis}.
\end{corollary}
\begin{proof}
The correspondence with encoders and Proposition~\ref{prop:score} prove equivalence. Existence and positivity follow from Theorem~\ref{thm:sdp}. Directly, the independent vectors $w_-,w_+$ bound $X\suc0$ on every finite objective sublevel; a peak block at one fixed input then bounds $a$. Closedness and the preceding nonaffineness argument ensure that every sublevel limit still has $X\succ0$.
\end{proof}

\section{Quartic Design}\label{app:quartic-proof}

\begin{proof}[Proof of Theorem~\ref{thm:quartic}]
Theorem~\ref{thm:rigidity} forces the quartic form for every exact encoder. Averaging the peak constraints at $s$ and $-s$ gives $A s^2+|c+v s^4|^2\le P$. Project $c$ onto the line of $v$ and choose the linear coefficient perpendicular to $v$. This preserves $A,B$, reduces the averaged norms, and removes $C$; Proposition~\ref{prop:score} shows that the score cannot decrease. Normalize $P=1$, write the resulting orthogonal curve as $\sqrt A\,s+i(bs^4+c)$ with $A,b>0$, and minimize the convex function
$$
 J(A,b,c)=M^2/A+1/b^2
$$
under the convex constraints
$$
 A t+(bt^2+c)^2\le1\qquad(0\le t\le1),
 \quad t=s^2.
$$
The proposed coefficients, with $A=a_\rho^2$ and $B=b_\rho^2$, obey the exact identity
\begin{equation}
 1-\left[A t+(b_\rho t^2+c_\rho)^2\right]
 =B(1-t)(t-\rho)^2(t+1+2\rho).
 \label{eq:quarticpeak}
\end{equation}
This proves every peak inequality. Equality holds at $t=\rho$ and $t=1$.

For global optimality, retain just these two active convex constraints and set
$$
 \eta=\frac{2}{B^2(1+\rho)^3(1-\rho)},\qquad
 \zeta=\eta\frac{1+\rho}{1-3\rho}.
$$
Both multipliers are positive. Let $y_\rho=b_\rho\rho^2+c_\rho=-b_\rho(1+\rho)^2/2$ and $y_1=b_\rho+c_\rho=b_\rho(1+\rho)(1-3\rho)/2$. The three stationarity identities are
\begin{equation}
\begin{aligned}
 -M^2/A^2+\eta\rho+\zeta&=0,\\
 -2/b_\rho^3+2\eta y_\rho\rho^2+2\zeta y_1&=0,\\
 2\eta y_\rho+2\zeta y_1&=0.
\end{aligned}
\label{eq:kkt}
\end{equation}
The last two follow by substitution. The first reduces exactly to \eqref{eq:quarticroot}, using $A=2B\rho(1+\rho)^2$. The Lagrangian is convex and has zero gradient at the proposed point, which therefore minimizes it globally. Any point satisfying all peak constraints satisfies these two constraints and has objective at least this minimum. Since the proposed curve also satisfies the full continuum by \eqref{eq:quarticpeak}, it is globally optimal. Scaling restores $P$ and gives \eqref{eq:quarticscore}.
\end{proof}

\section{Frontier Proofs}\label{app:frontier-proofs}

\subsection{General Frontier}\label{app:general-frontier-proof}
\begin{proof}
Corollary~\ref{cor:inverse-synthesis} identifies $\mathcal C_P$ with the continuous exact peak-admissible encoders and ensures $X\succ0$. Proposition~\ref{prop:score} gives the two endpoint budgets, and Proposition~\ref{prop:general-small-noise} gives the linear objective. Thus \eqref{eq:pareto-sdp} optimizes over the full continuous exact class.

A separation-optimal encoder supplies a feasible point for every stated floor. Since $Q_\pi\succ0$, a finite objective sublevel bounds $\operatorname{tr}X$ and hence $X$. The block constraints imply $|g(s)-a|^2\le P\operatorname{tr}X$, so they also bound $a$. The feasible sublevel is closed and bounded; every limit still has positive-definite $X$ by the preceding nonaffineness argument. Hence an optimum exists, including at zero floor.

To prove uniqueness and convexity in the floor, use the equivalent Gram description from Theorem~\ref{thm:sdp}. Feasible pairs $(H,\gamma)$ satisfy affine peak inequalities and $G-\gamma w_\sigma w_\sigma^\T\suc0$. Reducing the translation diagonal of $H$ to its Schur complement preserves $G$ and the floor and gives a planar encoder. The objective in these variables is $\operatorname{tr}(Q_\pi G^{-1})/(2K^2)$. It is strictly convex in $G$: its second directional derivative, before the positive scale factor, along a nonzero symmetric $D$ is
$$2\operatorname{tr}(A B^2)>0,$$
where $A=G^{-1/2}Q_\pi G^{-1/2}\succ0$ and $B=G^{-1/2}DG^{-1/2}\ne0$.
Convex combinations in the Gram formulation therefore prove convexity of $\mathcal V_\pi$ and uniqueness of its optimal $G$, equivalently $X$. Nested feasible sets prove monotonicity. Compact objective sublevels give lower semicontinuity; together with finiteness, convexity, and monotonicity this also gives continuity at both endpoints.

The unconstrained minimizer remains feasible precisely through its actual floor $\gamma_c$, proving the constant portion. Above $\gamma_c$, an optimal Gram matrix cannot have actual floor strictly greater than the required $\gamma$: mixing its feasible $H$ slightly with one realizing $X_0^{-1}$ preserves that floor and strictly lowers the objective. Thus the floor is active. Equal optimal values at two distinct floors above $\gamma_c$ would, by uniqueness at the smaller floor, give the same Gram matrix, contradicting activity. The same argument excludes equality between the value at $\gamma_c$ and any larger floor.
\end{proof}

\subsection{Quadratic Frontier}\label{app:quadratic-frontier-proof}
\begin{proof}
Write $G=\left(\begin{smallmatrix}A&C\\C&B\end{smallmatrix}\right)$. The reflection and translation-projection construction used for Theorem~\ref{thm:quartic}, applied here to $s^2$, preserves feasibility and $A,B$ while setting $C=0$. This decreases
$$
 \operatorname{tr}(Q_\pi G^{-1})
 =\frac{M_{\rm eff}^2B+A}{AB-C^2}
$$
and increases the worst-case score, so it preserves every robustness floor. Normalize $P=1$ and write the resulting curve as $\sqrt A\,s+i(bs^2+c)$, with $b=\sqrt B>0$. Its energy $At+(bt+c)^2$ is convex in $t=s^2\in[0,1]$, so the peak constraint is equivalent to
$$c^2\le1,\qquad A+(b+c)^2\le1.$$
If $b\le1$, these inequalities give $A\le1$ and $B\le1$; the feasible choice $A=B=1$ cannot worsen either objective. If $b>1$, $|c|\le1$ gives $A\le1-(b-1)^2=b(2-b)$, with equality attained at $c=-1$. Positive $A$ then requires $b<2$. Thus the one-parameter family \eqref{eq:quadratic-pareto-map} contains a global optimizer for the present constrained problem, not merely an optimum within a stipulated family.

The function $x^2/[b(2-b)]+1/b^2$ strictly decreases and then strictly increases, with its unique minimum at $b(x)$. Its stationarity equation is $x^2b(b-1)=(2-b)^2$, which also shows that $b(x)$ decreases with $x$. Hence $b_{\rm eff}>b_M$. The feasible set $R_M(b)\ge\gamma/P$ is an interval containing $b_M$, with right endpoint $b_+(\gamma/P)$. Minimizing the prior objective over that interval gives the stated clipped choice. Scaling restores $P$.
\end{proof}

\section{Grid Certificate}\label{app:grid-proof}
The grid problem gives an upper bound by omitting peak constraints. To obtain a feasible continuous encoder, we bound the energy excess between samples and remove it by scaling. Fixed spanning inputs make this bound uniform as the grid is refined.
\begin{proof}[Proof of Theorem~\ref{thm:certifiedgrid}]
Differentiating the peak function $q_H$ defined in Section~\ref{sec:certified} gives
$$
 |q_H''(s)|\le
 2\|H\|_2\bigl(\|z'(s)\|_2^2+
 \|z(s)\|_2\|z''(s)\|_2\bigr).
$$
On each interval of length at most $h$, the linear-interpolation remainder is bounded in absolute value by $M_hh^2/8$. The sampled peak inequalities therefore imply $q_{H_h}(s)\le P+M_hh^2/8$ throughout $I$. Scaling preserves the homogeneous semidefinite constraints and enforces the continuous peak bound. Dropping peak constraints gives the upper inequality in \eqref{eq:certifiedgrid}.

Let $Z\in\R^{3\times3}$ have columns equal to the three fixed spanning vectors $z(s)$. As in the proof of Theorem~\ref{thm:sdp}, every grid-feasible Gram matrix satisfies $\|H\|_2\le3P\|Z^{-1}\|_2^2$. This uniform bound proves existence of a grid optimizer, bounds $M_h$ independently of $h$, and bounds $U_h$ through $U_h\le G_{22}$. Consequently, $U_h(1-\rho_h)=O(h^2)$.
\end{proof}

\section{Proof of the Unified Extension}\label{app:polynomial}

We prove Theorem~\ref{thm:extension}, verify the stated sufficient conditions, and derive its quadratic specialization. Throughout, the two preprocessing functions and $F$ are prescribed before encoder design.

The representation $f=F\circ\boldsymbol\mu_\gamma$ also describes all continuous functions that admit identifiable computation with a shared continuous planar encoder when the preprocessing functions may be chosen freely with only continuity imposed on $\gamma$. Indeed, any encoder satisfying the identifiability condition supplies $\gamma=\enc$; this condition makes $F(\boldsymbol\mu_\gamma(\mathbf s))=f(\mathbf s)$ well defined, and the compact-quotient argument following \eqref{eq:exactness} makes $F$ continuous. Conversely, continuity and boundedness of $\gamma$ permit a continuous encoder $\enc=\rho\gamma$ with $\rho>0$ satisfying any prescribed $P>0$, after which the receiver recovers $\boldsymbol\mu_\gamma$ by division by $K\rho$.

\subsection{Classification Under the Identifiability Condition}

\begin{proof}[Proof of the classification in Theorem~\ref{thm:extension}]
Linear independence of $1,\varphi_1,\varphi_2$ implies that two interior tangents of $\gamma$ are independent. Otherwise all tangent vectors span at most a line, so some nonzero $a\in\R^2$ satisfies $a^\T\gamma'=0$ on $J$. Integration and endpoint continuity would make $a^\T\gamma$ constant on $I$, contradicting independence. Fix an independent tangent pair $\gamma'(y_*),\gamma'(z_*)$.

\emph{Local factorization.} If two encoded triple sums agree, appending the same fourth input $x\in J$ and the same $K-4$ fillers gives equal desired-function values $R_x$ for every $x$. When both triple statistics lie in one identifying neighborhood $B$ from \eqref{eq:completion-id}, their equality follows. Thus, locally in $B$, an encoded triple sum determines its two-dimensional statistic $p$.

Fix any first input $x_0\in J$. Near $(y_*,z_*)$, the pair map $(y,z)\mapsto T_3(x_0,y,z)$ is a local $C^1$ diffeomorphism. Its image is open in $\Omega_3$ and therefore meets the dense identifying set $\mathcal O$. Choose $y_0,z_0$ with independent tangents and $p_0=T_3(x_0,y_0,z_0)\in\mathcal O$. The inverse function theorem supplies coordinates
$$
 (x,y,z)\longmapsto(p,x)=(T_3(x,y,z),x)
$$
on a product $B\times L$ about $(p_0,x_0)$, with $B$ contained in an identifying neighborhood. Write the encoded triple sum as $A(p,x)$. The preceding implication makes $q(A(p,x))=p$ well defined on the image. Lemma~\ref{lem:topology} makes $A$ independent of $x$. Hence a continuous $H$ on $B$ satisfies, on a sufficiently small source box,
\begin{equation}
 \enc(x)+\enc(y)+\enc(z)
 =H\bigl(\gamma(x)+\gamma(y)+\gamma(z)\bigr).
 \label{eq:general-localfactor}
\end{equation}

\emph{Automatic differentiability.} The pair map $(y,z)\mapsto w=\gamma(y)+\gamma(z)$ has a local $C^1$ inverse $(y(w),z(w))$ near $(y_0,z_0)$. Choose a smooth density $\rho$ of integral one with compact support inside its image. Integrating \eqref{eq:general-localfactor} gives
$$
 \enc(x)=\int H(\gamma(x)+w)\rho(w)\,dw-C_0,
$$
where $C_0=\int[\enc(y(w))+\enc(z(w))]\rho(w)\,dw$ is independent of $x$. Shrink the source interval and the support of $\rho$ so that all arguments of $H$ lie in a compact subset of $B$. Multiplying $H$ by a smooth cutoff equal to one near that set and extending by zero gives a compactly supported continuous $\widetilde H$. The function
$$
 \mathcal B(a)=\int\widetilde H(a+w)\rho(w)\,dw
             =\int\widetilde H(z)\rho(z-a)\,dz
$$
is smooth: derivatives fall on $\rho$, and compact support justifies differentiation under the integral. Therefore $\enc=\mathcal B\circ\gamma-C_0$ is $C^1$ near $x_0$. Since $x_0$ was arbitrary, $\enc\in C^1(J;\R^2)$.

\emph{Constant affine coefficients.} At a regular triple whose statistic lies in $\mathcal O$, permute the inputs if necessary to select an independent tangent pair in the local factorization above. For a real component $e$ of $\enc$, differentiating along a local fiber of $T_3$ then gives
$$
 (e'(x),e'(y),e'(z))\in\operatorname{rowspan}DT_3(x,y,z).
$$
The implicit function theorem identifies the fiber tangent space with $\ker DT_3$, and the derivative of the encoded sum annihilates this kernel. Near every regular triple, $T_3$ maps open neighborhoods to open sets. Density of $\mathcal O$ therefore supplies regular triples with statistics in $\mathcal O$ approaching any given regular triple. Continuity of the derivative rows, and the vanishing of their $3\times3$ determinant, extend the row-space inclusion to every regular triple.

Fix $y=y_*$ and $z=z_*$. Their two entries determine unique constants $a_e,b_e$ satisfying
$$
 \begin{pmatrix}\varphi_1'(y_*)&\varphi_2'(y_*)\\
 \varphi_1'(z_*)&\varphi_2'(z_*)\end{pmatrix}
 \begin{pmatrix}a_e\\b_e\end{pmatrix}
 =\begin{pmatrix}e'(y_*)\\e'(z_*)\end{pmatrix}.
$$
Every triple $(x,y_*,z_*)$ is regular, so the same constants give $e'(x)=a_e\varphi_1'(x)+b_e\varphi_2'(x)$ for all $x\in J$. Integration and endpoint continuity yield $\enc=c+\mathsf U\gamma$ with constant $c\in\R^2$ and $\mathsf U\in\R^{2\times2}$.

If $\mathsf U$ were singular, an identifying open neighborhood $B$ would contain distinct $p,\widetilde p$ with $\mathsf Up=\mathsf U\widetilde p$. These statistics are attained by interior triples, whose encoded sums $3c+\mathsf Up$ agree. The identifiability condition for $f$ under every common completion and \eqref{eq:completion-id} force $p=\widetilde p$, a contradiction. Thus $\det\mathsf U\ne0$. Conversely,
$$
 r_\enc(\mathbf s)=Kc+K\mathsf U\boldsymbol\mu_\gamma(\mathbf s)
$$
allows recovery of $\boldsymbol\mu_\gamma$ and evaluation of $F$, proving sufficiency.
\end{proof}

\subsection{Convex Design and Attainment}

\begin{proof}[Proof of the design claim in Theorem~\ref{thm:extension}]
For $\enc=c+\mathsf U\gamma$ and $G=\mathsf U^\T\mathsf U\succ0$, cancellation of $c$ gives
\begin{equation}
 \Lambda_{\sq}(\enc;f)
 =K^2\inf_{\substack{\mu,\nu\in\mathcal M\\F(\mu)\ne F(\nu)}}
 \frac{(\mu-\nu)^\T G(\mu-\nu)}{|F(\mu)-F(\nu)|^2}.
 \label{eq:general-score-appendix}
\end{equation}
Every pair in $\mathcal M$ is attained by input vectors, so this is an equality of the original pairwise infima.

Let $F$ have a Lipschitz constant $L>0$ on $\mathcal M$. A sufficiently small $\rho>0$ makes $\enc=\rho\gamma$ peak admissible. Its coefficient Gram matrix is feasible in \eqref{eq:general-sdp} with $\lambda=\rho^2/L^2>0$. Every continuous encoder with positive score satisfies the identifiability condition and hence has the classified form. Its Gram matrix is feasible with $\lambda=\Lambda_{\sq}/K^2$ by \eqref{eq:general-score-appendix}. Thus the program bounds all positive-score competitors; zero-score encoders cannot improve the positive feasible value.

We next show why no rank constraint on the $3\times3$ matrix $H$ is needed. Let $(H,\lambda)$ be feasible with $\lambda>0$, and let $G$ be its lower-right block. If $G$ were singular, choose $0\ne v\in\ker G$ and distinct $p,\widetilde p=p+\varepsilon v$ inside an identifying neighborhood. For every common completion, their attained averages differ by $\varepsilon v/K$. The pair constraint gives $R_x(p)=R_x(\widetilde p)$ for every $x\in J$, contradicting \eqref{eq:completion-id}. Consequently $G\succ0$.

Write
$$
 H=\begin{pmatrix}h&\ell^\T\\\ell&G\end{pmatrix},\qquad
 \mathsf U^\T\mathsf U=G,\qquad c=\mathsf U^{-\T}\ell.
$$
Positive semidefiniteness gives $h\ge\ell^\T G^{-1}\ell=\|c\|^2$. The actual Gram matrix of the planar coefficient vectors $c$ and the columns of $\mathsf U$ therefore differs from $H$ only by reducing its top-left entry to $\|c\|^2$. This preserves $G$, reduces every peak quadratic form $z_\gamma(s)^\T Hz_\gamma(s)$ by the same nonnegative amount, and leaves all pair constraints unchanged. Hence $c+\mathsf U\gamma$ is a planar peak-admissible encoder with score at least $K^2\lambda$.

Finally, independence provides $s_1,s_2,s_3\in I$ for which $Z=[z_\gamma(s_1)\ z_\gamma(s_2)\ z_\gamma(s_3)]$ is invertible. The peak constraints and $H\succeq0$ imply
$$
 \operatorname{tr}(Z^\T HZ)\le3P,\qquad
 \|H\|_{\mathrm{op}}\le3P\|Z^{-1}\|_{\mathrm{op}}^2.
$$
Fix an attainable pair $\mu_0,\nu_0$ with unequal $F$ values. Its pair constraint gives
$$
 \lambda\le
 \frac{\|H\|_{\mathrm{op}}\|\mu_0-\nu_0\|^2}
 {|F(\mu_0)-F(\nu_0)|^2}.
$$
The feasible set is closed and bounded in a finite-dimensional space, so its positive optimum $\lambda_\gamma^\star$ is attained. The planar construction attains $K^2\lambda_\gamma^\star$, because a larger score would contradict the program's upper bound.
\end{proof}

\subsection{Sufficient Conditions for Identification}

For the original family, choose $\gamma(s)=(s,\phi(s))^\T$, $F(a,b)=b-\psi(a)$, and $s_\circ=0$. If two triple statistics $(S,T)$ and $(\widetilde S,\widetilde T)$ have equal values under every common completion, then
$$
 T-\widetilde T
 =K\left[\psi\!\left(\frac{S+x}{K}\right)
 -\psi\!\left(\frac{\widetilde S+x}{K}\right)\right]
 \qquad(x\in J).
$$
The strict-increment argument in Lemma~\ref{lem:collision}, comparing $x=0$ and $x=1/2$, gives $S=\widetilde S$ and then $T=\widetilde T$. Thus \eqref{eq:completion-id} holds globally on the regular triple-sum image. Nonaffineness of $\phi$ supplies independence of $1,s,\phi(s)$.

We next prove that \eqref{eq:analytic-nondegenerate} is sufficient for completion identifiability. Assume $F$ is real analytic on all of $\R^2$ and has no nonzero constant direction annihilating its gradient. Put
$$
 \widehat F(p)=F((p+d_\circ)/K),\qquad g=\nabla\widehat F.
$$
For $a,b\in J$, the determinant of the derivative of $p\mapsto(R_a(p),R_b(p))$ is
$$
 D_{a,b}(p)=\det[g(p+\gamma(a)),g(p+\gamma(b))].
$$
Suppose every such determinant vanished identically. Differentiating with respect to $b$ at $b=a$ and translating $p$ gives
$$
 \det[g(z),Dg(z)\gamma'(a)]=0
 \qquad(z\in\R^2,\ a\in J).
$$
Two independent tangents imply $\det(g,\partial_jg)=0$ for $j=1,2$. The nondegeneracy assumption excludes $g_1\equiv0$. On a connected open ball where $g_1\ne0$, both derivatives of $g_2/g_1$ vanish, so $g_2=\kappa g_1$ for a constant $\kappa$. The real-analytic identity theorem extends this equality to $\R^2$, making $(-\kappa,1)^\T g\equiv0$. Translation and scalar rescaling of the argument give the same forbidden direction for $\nabla F$.

Some fixed $a,b$ therefore have $D_{a,b}$ not identically zero. Its nonzero set intersected with $\Omega_3$ is dense and open. At every point of this set, the inverse function theorem makes $(R_a,R_b)$ locally injective, which implies \eqref{eq:completion-id}. Analyticity on the whole plane is used here; an analytic extension from a smaller domain is not assumed.

For polynomial $F$, failure of \eqref{eq:analytic-nondegenerate} is equivalent to $F(z)=h(w^\T z)$ for some nonzero $w$ and univariate polynomial $h$. Indeed, a constant direction annihilating $\nabla F$ makes $F$ constant on every line in that direction; orthogonal linear coordinates reduce $F$ to a polynomial in the remaining coordinate. Conversely, differentiating $h(w^\T z)$ supplies such an annihilating direction. This proves the polynomial criterion without restrictions on coefficient signs, parity, or degree.

\subsection{Quadratic Specialization and Examples}

Assume $\gamma(s)=(s,s^2)^\T$, completion identifiability, and the stated $C^1$ regularity and bounds \eqref{eq:general-endpoints} on $\mathcal C_2=\{(m,q):m^2\le q\le1\}$. For an encoder $\enc=c+us+vs^2$ satisfying the identifiability condition, write
$$
 G=\begin{pmatrix}A&C\\C&B\end{pmatrix}
 =\begin{pmatrix}\|u\|^2&u^\T v\\u^\T v&\|v\|^2\end{pmatrix}\succ0.
$$
Convexity of $w\mapsto w^\T G^{-1}w$ shows that its maximum on the gradient rectangle is attained at a corner. The endpoint condition gives
$$
 M_G=\max_{\sigma\in\{+,-\}}w_\sigma^\T G^{-1}w_\sigma
 =\frac{\alpha^2B+\beta^2A+2\alpha\beta|C|}{AB-C^2}.
$$
For any $\mu,\nu\in\mathcal M$, the joining segment lies in the convex set $\mathcal C_2$. Integrating the gradient on that segment and applying Cauchy--Schwarz in the metric $G$ gives
$$
 |F(\mu)-F(\nu)|\le\sqrt{M_G}\sqrt{(\mu-\nu)^\T G(\mu-\nu)}.
$$
Equation~\eqref{eq:general-score-appendix} therefore gives $\Lambda_{\sq}\ge K^2/M_G$.

To prove sharpness using attainable pairs, choose interior input vectors approaching the consensus endpoint whose gradient attains $M_G$, with at least two distinct components at each step. For example, near the positive endpoint use $(1-2\varepsilon,1-\varepsilon,\ldots,1-\varepsilon)$ with $\varepsilon\downarrow0$. The derivative of the moment map is
$$
 \frac1K\begin{pmatrix}1&\cdots&1\\2s_1&\cdots&2s_K\end{pmatrix},
$$
which has rank two at each such vector. An open neighborhood of each corresponding moment pair $\mu_n$ therefore lies in $\mathcal M$. Set $g_n=\nabla F(\mu_n)$, which is nonzero for large $n$. Small displacements $tG^{-1}g_n$ are attainable, and their separation quotients tend as $t\to0$ to $K^2/(g_n^\T G^{-1}g_n)$. Letting $n\to\infty$ gives the reverse bound. Thus
\begin{equation}
 \Lambda_{\sq}(\enc;f)
 =\frac{K^2(AB-C^2)}{\alpha^2B+\beta^2A+2\alpha\beta|C|}.
 \label{eq:polynomial-score}
\end{equation}
The convex hull supplies the lower bound, whereas actual finite-$K$ moment pairs supply the upper bound.

Put $M=\alpha/\beta$. The objective in \eqref{eq:polynomial-score} is $\beta^{-2}$ times the quadratic objective in Appendix~\ref{app:quadratic-proof}. Its peak-power certificate applies to arbitrary coefficient triples $(c,u,v)$ and uses only the scalar $M$. It gives
$$
 A=Pb(2-b),\qquad B=Pb^2,\qquad C=0,
$$
where $b\in(1,2)$ uniquely satisfies
$$
 \alpha^2b(b-1)=\beta^2(2-b)^2.
$$
Uniqueness follows because the left side increases strictly on $[1,2]$, while the right side decreases strictly. Solving for $b$ gives $b_{\alpha,\beta}$ and \eqref{eq:polyopt}. The proposed curve obeys
$$
 |\enc^\star(s)|^2/P=1-b^2s^2(1-s^2)\le1
 \qquad(s\in I).
$$
Encoders that fail the identifiability condition have zero score, so the positive optimum is global over all continuous peak-admissible encoders.

For the polynomial example \eqref{eq:polyclass}, nonnegative coefficients yield $|F_m|\le\alpha$ and $|F_q|\le\beta$ on $[-1,1]\times[0,1]$. Common parity $\epsilon$ of the powers of $m$ gives
$$
 w_+=(\alpha,\beta)^\T,\qquad
 w_-=((-1)^{\epsilon+1}\alpha,(-1)^\epsilon\beta)^\T,
$$
which supply the four corners together with their negatives. To verify identification, suppose $F(m,q)=h(am+bq)$. Dependence on both moments forces $a,b\ne0$ and $d=\deg h\ge1$. Its highest-total-degree part is a nonzero multiple of $(am+bq)^d$, containing nonzero terms with $m$ exponents zero and one. Lower-degree terms cannot cancel either term, contradicting common parity. The analytic criterion applies, proving the original polynomial result as an example of the extension.

Finally, $F(m,q)=1-e^{-(q-m^2)}$ is globally real analytic with
$$
 \nabla F(m,q)=e^{-(q-m^2)}(-2m,1)^\T.
$$
On $\mathcal C_2$, $0\le q-m^2\le1$, so the gradient bounds hold with $(\alpha,\beta)=(2,1)$. At the consensus endpoints, the gradients are $(-2,1)^\T$ and $(2,1)^\T$, which also exclude any constant annihilating direction. The analytic criterion and quadratic specialization apply. This nonpolynomial desired function therefore has the same optimal encoder and squared-separation score as variance.

An affine output change $\widetilde f=\eta f+\delta$, with $\eta,\delta\in\R$ and $\eta\ne0$, preserves the optimal encoder and divides its score by $\eta^2$. A fixed interval homeomorphism $h:I\to I$ likewise gives the encoder $\enc^\star\circ h$ for $f(h(s_1),\ldots,h(s_K))$, with the same peak constraint and score.

\section{Scope of the Criterion}\label{app:boundaries}
\subsection{Other Dissimilarities}
The shrinking-pair statement gives a criterion-transfer corollary over the full continuous competitor class.

\begin{corollary}
Suppose $D:[0,\infty)\to[0,\infty)$ is continuous, $D(0)=0$, positive away from zero, and
$$
 D(d)\ge c_Dd^2\quad(0\le d\le2K\sqrt P),\qquad
 \lim_{d\downarrow0}\frac{D(d)}{d^2}=c_D>0.
$$
For the targets of Proposition~\ref{prop:score}, every exact continuous encoder satisfies $\Lambda_D=c_D\Lambda_{\sq}$. Therefore every optimal encoder for the targets of Proposition~\ref{prop:score} remains optimal, and its optimum is multiplied by $c_D$.
\end{corollary}
\begin{proof}
The lower envelope gives $\Lambda_D\ge c_D\Lambda_{\sq}$. Proposition~\ref{prop:score} supplies pairs attaining the opposite inequality in the limit. Nonexact encoders have zero score because $D(0)=0$.
\end{proof}

Examples are $D(d)=d^2+ad^4$ with $a\ge0$ and $D(d)=e^{ad^2}-1$ with $a>0$. The original uncentered $e^{ad^2}$ criterion has $D(0)=1$ and is not covered. No closed-form optimizer for that criterion is claimed here. Gaussian MSE and optimality at a fixed disturbance radius also remain separate problems.


\subsection{Exactness and Conditioning}
The family
$$\enc_\beta(s)=\sqrt P[s+i\beta(1-s^2)],\qquad0<\beta\le1,$$
computes $mq$ exactly, but its squared score is
$$\Lambda_\sq(\enc_\beta;mq)=K^2P\frac{\beta^2}{1+\beta^2}.$$
Its inverse conditioning therefore diverges as $\beta\downarrow0$. At zero and $K=4$, inputs $(1,1,0,0)$ and $(1/2,1/2,1/2,1/2)$ give the same received sum while their targets are $1/4$ and $1/8$.

For a nonconstant bounded target, the uncentered exponential criterion assigns every encoder a positive score of at least $1/(\max f-\min f)^2$, even if exactness fails. Its nonzero value at zero distance therefore makes exactness a separate condition. Noise-dependent exponential optimization requires a separate analysis.

\subsection{Finite Alphabets}
For $K=4$ and $s_k\in\{-1,1\}$, the possible means are $m\in\{-1,-1/2,0,1/2,1\}$ and variance is $1-m^2$. Antipodal signaling gives received sum $4\sqrt P m$, hence
$$\frac{|\Delta r|^2}{|\Delta V|^2}=\frac{16P}{(m+\widetilde m)^2}$$
whenever the target gap is nonzero. The largest admissible $|m+\widetilde m|$ is $3/2$, giving score $64P/9$. Every binary map has symbol separation at most $2\sqrt P$, so antipodal signaling maximizes every received gap and is globally optimal. Restricting the continuous variance-optimal parabola instead gives
$$\left(\frac4{\sqrt3}-\frac43\right)\frac{64P}{9}<\frac{64P}{9}.$$
There is no contradiction: removing source values removes the continuum's sharp shrinking pairs.


\section{Approximate Rigidity}
\label{app:approx-rigidity}
The encoder characterization also gives a qualitative statement for approximate implementations when the admissible encoder class is compact. The common regularity and decoder bounds below are explicit additional assumptions.

Fix the desired function $f_{\phi,\psi}:I^K\to\R$ and impose the hypotheses of Theorem~\ref{thm:rigidity}. Let $\omega:[0,2]\to[0,\infty)$ satisfy $\omega(h)\to0$ as $h\downarrow0$. Let $\mathcal E_{P,\omega}$ consist of the encoders $E\in C(I;\R^2)$ satisfying
$$
 \|E\|_\infty\le\sqrt P,\quad
 |E(s)-E(t)|\le\omega(|s-t|)\quad\forall s,t\in I.
$$
Write $\mathcal A_\phi\subset C(I;\R^2)$ for the family $s\mapsto c+us+v\phi(s)$ with $c,u,v\in\R^2$ and linearly independent $u,v$, and use uniform distance to this family.

\begin{proposition}
\label{prop:approx-rigidity}
For fixed $P$, $\omega$, and $L>0$, there is a nondecreasing function $\eta:[0,\infty)\to[0,\infty)$ with $\lim_{\epsilon\downarrow0}\eta(\epsilon)=0$ and the following property. For every $\epsilon\ge0$, if $E\in\mathcal E_{P,\omega}$ and an $L$-Lipschitz decoder $D:\R^2\to\R$ satisfy
$$
 \sup_{\mathbf s\in I^K}
 |D(r_E(\mathbf s))-f_{\phi,\psi}(\mathbf s)|\le\epsilon,
$$
then $\operatorname{dist}_\infty(E,\mathcal A_\phi)\le\eta(\epsilon)$. The function $\eta$ may also depend on the fixed $K$, $\phi$, and $\psi$.
\end{proposition}
\begin{proof}
Approximate computation and decoder Lipschitz continuity imply, for every $\mathbf s,\mathbf t\in I^K$,
\begin{equation}
 |f(\mathbf s)-f(\mathbf t)|
 \le 2\epsilon+L|r_E(\mathbf s)-r_E(\mathbf t)|.
 \label{eq:approx-collision}
\end{equation}
The peak and common modulus bounds make $\mathcal E_{P,\omega}$ compact in the uniform norm by the Arzel\`a--Ascoli theorem. The class is closed because its defining inequalities survive uniform limits.

Define the continuous defect
$$
 d_L(E)=\frac12\sup_{\mathbf s,\mathbf t\in I^K}
 \bigl[|f(\mathbf s)-f(\mathbf t)|
       -L|r_E(\mathbf s)-r_E(\mathbf t)|\bigr]_+.
$$
Indeed, $|d_L(E)-d_L(\widetilde E)|\le KL\|E-\widetilde E\|_\infty$. Set $\eta(\epsilon)$ to the supremum of $\operatorname{dist}_\infty(E,\mathcal A_\phi)$ over $E\in\mathcal E_{P,\omega}$ with $d_L(E)\le\epsilon$, taking the supremum to be zero when the set is empty. The distance is continuous even though independence makes $\mathcal A_\phi$ nonclosed. Compactness makes this supremum finite, and \eqref{eq:approx-collision} proves the required bound.

If $\eta(\epsilon)$ failed to tend to zero, a sequence $\epsilon_n\downarrow0$ would supply encoders at a fixed positive distance from $\mathcal A_\phi$. A subsequence converges uniformly to $E\in\mathcal E_{P,\omega}$. Passing to the limit gives
$$
 |f(\mathbf s)-f(\mathbf t)|
 \le L|r_E(\mathbf s)-r_E(\mathbf t)|.
$$
Thus $f$ is identifiable through $E$, and Theorem~\ref{thm:rigidity} places $E$ in $\mathcal A_\phi$ with linearly independent coefficient vectors $u,v\in\R^2$. Uniform convergence then contradicts the positive distance. This proof requires no convergence of the decoders.
\end{proof}

This proposition supplies neither an explicit rate nor a quantitative parameter-identification guarantee. Its assumptions hold, for example, for encoders with a prescribed common Lipschitz constant and decoders with a fixed disturbance-amplification budget. Peak power alone does not supply the compactness used in the proof. We do not assert that these additional assumptions are necessary or that the proposition resolves quantitative approximate rigidity.

\section{Lean Verification}\label{app:lean}

We verify selected results in Lean~4~\cite{lean4}, using mathlib~\cite{mathlib}. The artifact includes the proof sources, a theorem map, a complete manuscript coverage inventory, and reproducible compilation and axiom audits. It distinguishes complete results from specializations and conditional steps.

The general collision lemma is verified for arbitrary features and encoders and every $K\ge4$, without differentiability of $\psi$. The operational score formula \eqref{eq:exactscore} and its shrinking-pair clause are also verified for every nondegenerate two-feature encoder and $K\ge2$, under the stated feature curvature and endpoint regularity assumptions. The formal score is the infimum over actual admissible input pairs. The proof uses explicit two-input perturbations and one-sided endpoint derivatives.

For variance, the score identity additionally covers degenerate quadratic curves. Lean verifies the operational optimum over every quadratic coefficient curve at all positive powers:
$$
 \max_{\substack{\enc(s)=c+us+vs^2\\|\enc(s)|^2\le P\;\forall s\in I}}
 \Lambda_{\sq}(\enc;V)
 =\frac{4K^2P}{9+6\sqrt3},\qquad K\ge2.
$$
The complete quadratic and quartic Gram-objective optimizations, including continuum peak certificates and arbitrary positive slope parameter and power, are also checked. These coefficient-class theorems do not yet extend their competitors to every continuous encoder.

For encoder design, the artifact verifies the Schur-complement reduction, rank-one score constraints, and explicit planar reconstruction. The reconstruction is connected to the general operational score: a positive feasible SDP point yields a continuous, exact encoder satisfying the full peak constraint and the promised separation bound. The proof derives spanning anchors and a positive feasible point from the feature assumptions, establishes compactness, and proves maximum-separation attainment and its operational realization within the affine-feature class.

The rigidity development constructs genuine local moment inverses and neighborhoods, derives the derivative-kernel relations from local fiber constancy, integrates through both source endpoints, and proves coefficient independence from exactness. Its final representation theorem still assumes interior encoder differentiability and local moment-fiber constancy. Deriving these properties from continuity and exactness requires the remaining topological factorization and smoothing arguments; planar invariance of domain is not formalized here.

The decoder development verifies the sharp scalar amplification principle, including the zero-score impossibility result, and the explicit complex-channel variance decoder with corrected strip projection. It also checks coherent implementation-error bounds and peak-safe contraction before quantization. The small-noise development proves the derivative limit, integrable domination, dominated convergence in noise-power coordinates, and independent isotropic second-moment evaluation. It derives local differentiability of the actual variance decoder and its leading-MSE limit. The complete general-decoder gradient-to-trace specialization in Proposition~\ref{prop:general-small-noise} remains unfinished.

For the probabilistic design problem, Lean verifies convexity of the continuum Gram constraints, strict convexity of inverse-Gram risk, and convexity and monotonicity of the actual infimum frontier. The artifact also proves existence of a risk-minimizing Gram design at every floor from zero through the positive maximum, supplying the optimizer witnesses used by its uniqueness, plateau, activity, and strict-increase results. Endpoint continuity, the five-variable inverse-form equivalence, and the complete probabilistic and continuous-encoder links remain unfinished.

The artifact pins Lean~4.24.0 and a matching mathlib revision. Its audited declarations use only \texttt{propext}, \texttt{Classical.choice}, and \texttt{Quot.sound}, with no unfinished proofs or additional axioms. The unrestricted continuous-encoder theorem, complete probabilistic frontier, certified-grid analysis, and remaining appendix results are not fully formalized. The artifact therefore supports a claim about selected verified results, not all proofs in the paper.

\section*{Acknowledgment}
The author acknowledges the use of ChatGPT (OpenAI) for assistance with language editing and coding\ifArxivVersion, including formalizing proofs in Lean\fi. All technical results and recorded experiments are entirely the author's original work.

% \ifdefined\SubmissionFormat\else\balance\fi
% \input{bibliography}
\bibliographystyle{IEEEtran}
\bibliography{bibliography}

\end{document}
